\ifdefined\pdfsuppressptexinfo
  \pdfsuppressptexinfo=15
\fi
\documentclass[11pt]{article}
\pdfinfoomitdate=1
\pdftrailerid{}

\usepackage[T1]{fontenc}
\usepackage[utf8]{inputenc}
\usepackage{lmodern}
\usepackage[margin=1in]{geometry}
\usepackage{amsmath,amssymb,amsthm,mathtools}
\usepackage{booktabs}
\usepackage{array}
\usepackage{longtable}
\usepackage{enumitem}
\usepackage{microtype}
\usepackage{xcolor}
\usepackage{hyperref}
\makeatletter\g@addto@macro{\UrlBreaks}{\do\_}\makeatother

\hypersetup{
  colorlinks=true,
  linkcolor=blue!55!black,
  citecolor=blue!55!black,
  urlcolor=blue!55!black,
  pdftitle={Silver-rank zeros of the tensor-curl eta invariant on lens spaces: silver S-units, a Pell ladder and a period law},
  pdfauthor={Luqman Mahmood}
}

\newtheorem{theorem}{Theorem}[section]
\newtheorem{proposition}[theorem]{Proposition}
\newtheorem{lemma}[theorem]{Lemma}
\newtheorem{corollary}[theorem]{Corollary}
\newtheorem{conjecture}[theorem]{Conjecture}
\theoremstyle{definition}
\newtheorem{definition}[theorem]{Definition}
\newtheorem{remark}[theorem]{Remark}
\newtheorem{example}[theorem]{Example}
\newtheorem{observation}[theorem]{Observation}

\DeclareMathOperator{\sgn}{sgn}
\newcommand{\Z}{\mathbb Z}
\newcommand{\Q}{\mathbb Q}
\newcommand{\R}{\mathbb R}
\newcommand{\C}{\mathbb C}
\newcommand{\ind}{\mathbf 1}
\newcommand{\e}{\mathrm e}
\newcommand{\sawtooth}[1]{\left(\!\left(#1\right)\!\right)}
\newcolumntype{L}[1]{>{\raggedright\arraybackslash}p{#1}}

\title{Silver-rank zeros of the tensor-curl eta invariant\\
on lens spaces:\\
silver $S$-units, a Pell ladder, and a period law}
\author{Luqman Mahmood}
\date{1 October 2026\\[3pt]\small Preprint\\
DOI: \href{https://doi.org/10.5281/zenodo.23049615}{10.5281/zenodo.23049615}}

\begin{document}

\maketitle

\begin{abstract}
Let $\eta_J(L(p,q))$ be the eta invariant of the rank-$J$ projected
transverse-traceless tensor-curl operator on the round lens space
$L(p,q)$ \cite{MahmoodP1Collection2026}.  At rank one it is $-4s(q,p)$,
which vanishes exactly on the amphicheiral spaces $q^2\equiv-1\pmod p$,
so there a zero reads as mirror symmetry.  At higher rank this fails on
the silver levels $p\mid2J^2-1$, where $\eta_J$ is an even integer and
zeros of chiral spaces are abundant.  We study the zero set at the level
$p=2J^2-1$.  Reading $\Z[\sqrt2]$ modulo $p$ through $\sqrt2\mapsto2J$,
the four zero families of \cite{MahmoodP1Collection2026} become the
orbits of the silver $S$-units $\sqrt2$, $\delta/\sqrt2$, $\delta$ and
$2\delta$, $\delta=1+\sqrt2$.  We prove a fifth family,
$\eta_J(L(p,3-4J))=J-1$ for odd $J$ and $0$ for even $J$, and an infinite
Pell ladder on the orbits of the powers $\delta^n$: with $P_n$ the Pell
numbers, $\eta_J=\tfrac23P_{n-2}J$ for even $n$ and $0$ for odd $n$ when
$P_n\mid J$, and, for $n\geq2$, $\eta_J=\gamma_n(J-1)/P_n$ with an
explicit $\gamma_n>0$ when $J\equiv1\pmod{P_n}$.
We conjecture a period law (every nonconstant $S$-unit orbit is eventually
quasi-polynomial in $J$ with period dividing its coefficient of $\sqrt2$),
support it by a reported out-of-sample test on $64$ orbits, and prove it
for every nonconstant orbit of norm $\pm1$ or $\pm2$
(the powers $\delta^n$, $n\geq1$, and the norm-two orbits; the constant
orbit $q=\pm1$ lies outside the law and is cubic in $J$).  There, on
each residue class of $J$, $\eta_J$ is a linear function of $J$ past an
explicit threshold; the zero set of each orbit is computable by
finitely many exact evaluations and is eventually periodic, with density
equal to the proportion of zero classes.  Among these orbits exactly the
first three families of \cite{MahmoodP1Collection2026} vanish at every
rank, and in a window of $64$ orbits exactly the four families do.
Consequently, for every $K$ there is a rank at which $K$ pairwise
non-homeomorphic chiral lens spaces, none balanced level by level, all
have $\eta_J=0$.  The period law beyond norm $\pm2$, the zeros outside
the $S$-units and the complete zero set remain open.
\end{abstract}

\section{Introduction}
\label{sec:intro}

The paper \cite{MahmoodP1Collection2026} evaluates the eta invariant
$\eta_J(L(p,q))$ of the rank-$J$ projected transverse-traceless (TT)
tensor-curl operator on every round lens space and exhibits four
infinite families of zeros at the level $p=2J^2-1$:
$q=J$, $J+1$, $2J-1$ and $J(J-1)$.  The eta invariant measures the net
helicity of the spectrum, and a natural first guess is that it vanishes
only for a reason of symmetry.  At rank one that guess is right, by
composition of classical results: $\eta_1(L(p,q))=-4s(q,p)$
\cite{MahmoodP1Collection2026}, the Dedekind sum $s(q,p)$ vanishes iff
$q^2\equiv-1\pmod p$ \cite[p.~36]{RademacherGrosswald1972}, and
$q^2\equiv-1$ is the criterion for $L(p,q)$ to be amphicheiral
\cite{PrzytyckiYasuhara2003,Boldt2017}; see
\cite[Cor.~4.4]{LensChirality2026}.  The companion paper
\cite{LensChirality2026} shows that at prime moduli $p\geq5$ the same
dichotomy persists at every rank $J$ with $2J^2\not\equiv1\pmod p$ and
$p\nmid J$ (Corollary~\ref{cor:E}(c) below).  The exceptions are the
\emph{silver ranks} $2J^2\equiv1\pmod p$.  There $\eta_J$ is an even
integer, and zeros of chiral spaces are common: at the $31$ silver primes
in $[211,599]$ about $15\%$ of the $q$-orbits are zeros
\cite[Obs.~5.6]{LensChirality2026}.  This paper is about those zeros, at
the level $p=2J^2-1$ itself.

\paragraph{Results.}
\begin{enumerate}[label=(\arabic*)]
\item \emph{A coordinate.}  The ring homomorphism
$\Z[\sqrt2]\to\Z/p$, $\sqrt2\mapsto2J$, sends the silver $S$-units
$\pm\delta^n\sqrt2^{\,m}$ ($\delta=1+\sqrt2$) to units, and the four
families of \cite{MahmoodP1Collection2026} are the orbits of $\sqrt2$,
$\delta/\sqrt2$, $\delta$ and $2\delta$ (Lemma~\ref{lem:sunit-four}).
The Cayley coordinate $\rho=(1+q)/(1-q)$ gives a dictionary for the
known zero orbits (Proposition~\ref{prop:cayley-dictionary}).
\item \emph{A fifth family} (Theorem~\ref{thm:C4}):
$\eta_J(L(2J^2-1,\,3-4J))=(J-1)\ind_{J\text{ odd}}$, the orbit of
$\delta^{-2}$.
\item \emph{The Pell ladder} (Theorem~\ref{thm:P}): exact values of
$\eta_J$ on the orbit of $\delta^n$ on the two slices $P_n\mid J$ and
$J\equiv1\pmod{P_n}$.  The first two rungs are the third family of
\cite{MahmoodP1Collection2026} and the fifth family; the orbit of
$\delta^n$ vanishes at every rank only for $n=1$.
\item \emph{The period law.}  Conjecture~\ref{conj:PL} with its out-of-sample
evidence, and Theorem~\ref{thm:PL}, which proves it for every nonconstant
orbit of norm $\pm1$ or $\pm2$ (the constant orbit $q=\pm1$ is outside
the law and cubic in $J$), with thresholds, linear class polynomials, a
decision procedure and complete zero sets for $24$ families, including one isolated zero
($J=429$) that is not part of any zero class.  The zero sets are
eventually periodic, with an explicit density
(Corollary~\ref{cor:eventual}).
\item \emph{Rank-uniform zeros.}  Among the orbits of norm $\pm1$ and
$\pm2$, exactly the first three families of
\cite{MahmoodP1Collection2026} vanish at every rank
(Corollary~\ref{cor:R}); in a window of $64$ orbits, exactly the four
families do (Corollary~\ref{cor:U}(a)).
\item \emph{Unbounded ambiguity} (Corollary~\ref{cor:ambiguity}): for
every $K$ there is a rank at which $K$ pairwise non-homeomorphic chiral
lens spaces, none balanced level by level, all have $\eta_J=0$.  This is
the precise sense in which the rank-one dictionary ``zero $=$ mirror''
fails at higher rank: a single vanishing eta invariant can hide an
unbounded amount of handedness.
\end{enumerate}
Along the way we record a refuted conjecture (Conjecture~\ref{conj:C5})
and a refuted ``zero-term principle'' (Remark~\ref{rem:zero-term}), and
explain why the conjugacy mechanism that proves the family $q=J$ at every
divisor of $2J^2-1$ in \cite{LensChirality2026} does not close the
classification (Remark~\ref{rem:kill}).

\paragraph{What is proved, what is observed.}  Every statement is
labelled as proved, conjectural or observed.  Observations are finite
computations whose ranges are stated.  The Faulhaber-sum evaluations and
the leading-coefficient table are certificate-assisted
(Remarks~\ref{rem:faulhaber-certificate} and \ref{rem:period-certificates},
checked by two independent computer algebra routes).  The finite
threshold and class-polynomial reconstruction also uses exact algebraic
checks, as specified in Theorem~\ref{thm:PL}.  Exact computations
are listed in Appendix~\ref{app:verification}; they test the identities
and conventions, and do not replace the proofs.  No physical claim is
made.  The methods are elementary (Dedekind reciprocity, Pell arithmetic,
lattice counting and an integrality argument), and Dedekind sums at
Pell-type arguments are classical territory; we have not made a
literature search for the individual evaluations and claim no priority
for them.  Eventual quasi-polynomial behaviour of parametric counting
functions is also classical in general; for context see the parametric
Presburger arithmetic of Bogart, Goodrick and Woods \cite{BGW2017}.  We
do not claim quasi-polynomiality as such.  What is specific here is the
explicit decomposition of $\eta_J$ into a Dedekind term, a residue term
and a lattice box, the cancellation that makes every class linear, the
explicit thresholds, and the resulting zero-set procedure in these
sectors.

\paragraph{Relation to the companion papers.}  We use the operator, the
normalization and the master formula of
\cite{MahmoodP1Collection2026}, and the following results of the
companion \cite{LensChirality2026}, restated in
Section~\ref{sec:conventions}: the classification of level-wise balance,
the Dedekind-sum facts, the universal congruence
$\eta_J-J(2J^2-1)\eta_1\in2\Z$, the exact denominator and the
prime-modulus dichotomy, the zero $\eta_J(L(p,J))=0$ at every divisor
$p>1$ of $2J^2-1$, and the Dedekind sum at the silver unit.  No result of
\cite{LensChirality2026} depends on the present paper; its Section~9 summarizes the present
results.

\paragraph{Structure of the proofs.}  Section~\ref{sec:sunit} introduces
the two coordinates.  The fifth-family and Pell-ladder proofs evaluate
the Dedekind, residue and lattice terms separately.
Section~\ref{sec:period-law} states the period-law conjecture and its
finite evidence.  Section~\ref{sec:period-law-proved} proves the law in
the norm-one and norm-two sectors, using an integrality argument and
linear breakpoint bounds, and derives the zero-set procedure and
unbounded ambiguity.

\section{Conventions and imported results}
\label{sec:conventions}

We keep the conventions of \cite[Sec.~2]{MahmoodP1Collection2026} and
\cite[Sec.~2]{LensChirality2026}.  For $p\geq2$ and
$q\in(\Z/p\Z)^\times$ the generator of $\Z_p$ acts on
$S^3\subset\C^2$ by $(z_1,z_2)\mapsto(\e^{2\pi i/p}z_1,\e^{2\pi iq/p}z_2)$,
and $L(p,q)$ is the quotient with its complex quotient orientation.  We
write $q^*$ for an inverse of $q$ modulo $p$, $\langle k\rangle_p$ for
the least non-negative residue, and $\tilde q$ for the odd lift of $q$
($\tilde q=q$ if $q$ is odd, $\tilde q=q+p$ if $q$ is even and $p$ odd).
Put
\begin{equation}
 A=\frac{1+\tilde q}{2},\qquad B=\frac{1-\tilde q}{2},\qquad A+B=1.
 \label{eq:AB}
\end{equation}
Let $W_a=\{-a,-a+2,\ldots,a\}$ be the weights of the $SU(2)$
representation $V_a$.  The lifted generator acts on the weight vector
$(w_L,w_R)$ of $V_a\otimes V_b$ by $\exp(2\pi i(Aw_L+Bw_R)/p)$, and for
$a\equiv b\pmod2$ the invariant count is
\[
 I_{a+1,b+1}(p,q)=\#\{(w_L,w_R)\in W_a\times W_b:\ Aw_L+Bw_R\equiv0\pmod p\}.
\]
The rank-$J$ operator has signed spectrum $+(n+J+1)$ on
$V_{n+2J}\otimes V_n$ and $-(n+J+1)$ on $V_n\otimes V_{n+2J}$, $n\geq0$
\cite[Thm.~3.1]{MahmoodP1Collection2026}.  With
$D_J(n)=I_{n+2J+1,n+1}-I_{n+1,n+2J+1}$ the level imbalance, $\eta_J$ is
the value at $s=0$ of the continuation of $\sum_{n\geq0}D_J(n)(n+J+1)^{-s}$
\cite[Prop.~4.1]{MahmoodP1Collection2026}.  We use the lens symmetries
\begin{equation}
 \eta_J(L(p,q^*))=\eta_J(L(p,q)),\qquad
 \eta_J(L(p,-q))=-\eta_J(L(p,q)),
 \label{eq:lens-symmetries}
\end{equation}
the Dedekind sum
$s(q,p)=\sum_{a=0}^{p-1}\sawtooth{a/p}\sawtooth{qa/p}$, the rank-one value
$\eta_1(L(p,q))=-4s(q,p)$ \cite[Cor.~5.8]{MahmoodP1Collection2026}, and
the master formula \cite[Thm.~5.1]{MahmoodP1Collection2026}
\begin{align}
 \eta_J(L(p,q))={}&-4J\,s(q,p)
 +\frac2p\sum_{a=1}^{J-1}(J-a)\bigl[\rho_p(aq)+\rho_p(aq^*)\bigr]\nonumber\\
 &-2\!\!\sum_{\substack{a,b\geq1\\a+b\leq J-1}}\!\!(J-a-b)
 \bigl[\delta_p(aq-b)-\delta_p(aq+b)\bigr]
 -L_J(p,q),
 \label{eq:master}
\end{align}
with $\rho_p(k)=0$ if $p\mid k$ and $p-2\langle k\rangle_p$ otherwise,
$\delta_p(k)=\ind_{p\mid k}$, and the low-mode term
$L_J(p,q)=\sum_{m=1}^{J-1}[I_{J-m,J+m}(p,q)-I_{J+m,J-m}(p,q)]$.
Equation~\eqref{eq:master} holds for every $p\geq2$, every unit $q$ and
every $J\geq1$.

\begin{definition}
\label{def:terms}
Let $q$ be a unit modulo $p$ and $J\geq1$.  $L(p,q)$ is \emph{rank-$J$
balanced} if $D_J(n)=0$ for every level $n\geq0$; it is \emph{mirror}
(amphicheiral) if $q^2\equiv-1\pmod p$ and \emph{chiral} otherwise
\cite{PrzytyckiYasuhara2003,Boldt2017}.  The rank $J$ is a \emph{silver
rank} for $p$ if $2J^2\equiv1\pmod p$, and a \emph{blind rank} if
$p\mid2J$.  The \emph{Clifford condition} is $\gcd(q\mp1,p)\leq2$.
\end{definition}

The following results of \cite{LensChirality2026} are used; we restate
them with their numbering there.

\begin{theorem}[{\cite[Thm.~4.1]{LensChirality2026}}]
\label{thm:classification}
Let $p\geq3$, $q$ a unit and $J\geq1$.  $L(p,q)$ is rank-$J$ balanced iff
\textup{(a)} $q^2\equiv-1\pmod p$; or \textup{(b)} $p\mid4J^2-2$ and
$q\equiv\pm(2J\pm1)\pmod p$; or \textup{(c)} $p\mid2J$ and the Clifford
condition holds.
\end{theorem}

\begin{lemma}[{\cite[Lem.~6.2]{LensChirality2026}}]
\label{lem:dedekind-facts}
For $p\geq2$ and $q$ a unit, $T:=12p\,s(q,p)\in2\Z$; in particular
$6p\,s(q,p)\in\Z$.
\end{lemma}

\begin{theorem}[{Universal congruence, \cite[Thm.~6.7]{LensChirality2026}}]
\label{thm:U}
For every $p\geq2$, unit $q$ and $J\geq1$,
$\eta_J(L(p,q))-J(2J^2-1)\,\eta_1(L(p,q))\in2\Z$.
\end{theorem}

\begin{theorem}[{\cite[Thm.~6.8(c)]{LensChirality2026}}]
\label{thm:denominator}
Write $p=2^a3^bm$ with $\gcd(m,6)=1$ and $C=J(2J^2-1)$.  Then
$\eta_J(L(p,q))\in2\Z$ iff $m\mid C(q^2+1)$, $(a<2$ or $2^a\mid J)$, and
$(b=0$ or $3^{b+1}\mid J)$.
\end{theorem}

\begin{corollary}[{\cite[Cor.~6.4]{LensChirality2026}}]
\label{cor:E}
Let $p\geq5$ be prime.  \textup{(b)} If $2J^2\equiv1\pmod p$ then
$\eta_J(L(p,q))\in2\Z$ for every unit $q$.  \textup{(c)} If
$2J^2\not\equiv1\pmod p$ and $p\nmid J$, then $\eta_J(L(p,q))=0$ iff
$q^2\equiv-1\pmod p$.
\end{corollary}

For $J\geq2$ and a divisor $p>1$ of $2J^2-1$ put
\begin{equation}
 \varepsilon:=4J+3,\qquad \varepsilon^2-6\varepsilon+1=8(2J^2-1)\equiv0\pmod p,
 \label{eq:silver-unit}
\end{equation}
the silver unit $3+2\sqrt2$ read modulo $p$ with $\sqrt2\equiv2J$.

\begin{theorem}[{\cite[Thm.~7.1]{LensChirality2026}}]
\label{thm:C3}
For every $J\geq2$ and every divisor $p>1$ of $2J^2-1$,
$\eta_J(L(p,J))=0$, and hence $\eta_J(L(p,q))=0$ for $q\equiv\pm J$ and
$q\equiv\pm2J\pmod p$.
\end{theorem}

\begin{corollary}[{\cite[Cor.~7.11]{LensChirality2026}}]
\label{cor:S}
Let $p$ be odd and $\varepsilon^2-6\varepsilon+1\equiv0\pmod p$.  Then
$s(\varepsilon,p)=-(p-1)/(2p)$ if $p\equiv1\pmod8$ and $1/(2p)$ if
$p\equiv7\pmod8$.
\end{corollary}

\section{Zeros at silver ranks: what is known}
\label{sec:overview}

\begin{itemize}
\item \textup{(Proved.)}  At a silver rank of a prime $p\geq5$, $\eta_J(L(p,q))$ is an
even integer for every unit $q$ (Corollary~\ref{cor:E}(b)); a silver rank
exists iff $p\equiv\pm1\pmod8$.  At composite $p\mid2J^2-1$ (odd, coprime
to $3$), $\eta_J\in2\Z$ iff $p\mid J(2J^2-1)(q^2+1)$, which holds
(Theorem~\ref{thm:denominator}), so $\eta_J$ is an even integer at
every silver rank of every modulus.
\item \textup{(Proved zero orbits.)}  The balanced family
$q\equiv\pm(2J\pm1)$ (Theorem~\ref{thm:classification}(b)); the orbit
$q\equiv\pm J,\pm2J$ (Theorem~\ref{thm:C3}); the mirror orbit when
$-1$ is a square modulo $p$; and, when $p=2J^2-1$ exactly, the orbits of
$q=J+1$ and $q=J(J-1)$ of \cite[Thm.~7.1]{MahmoodP1Collection2026}, the
fifth family of Theorem~\ref{thm:C4} (even $J$), and the Pell ladder of
Theorem~\ref{thm:P}.
\item \textup{(Observed.)}  \cite[Obs.~5.5 and 5.6]{LensChirality2026}: at the $31$ silver primes in $[211,599]$,
$471$ of $3{,}076$ $q$-orbits ($15.3\%$) are zeros, with the zero fraction
highest at small silver ranks ($33\%$ at $p=241$, $J=11$); at
$(97,7)$ and $(199,10)$ the cube roots of unity are zeros.  In the
divisor family of Theorem~\ref{thm:C3} with $J\leq59$, further zeros
occur at roughly $25$--$30\%$ of the rows examined in the supplement
\cite{LCrecord}.  These fractions are
descriptive statistics of the examined ranges, computed with different
orbit conventions, and are not asserted to have limits.
\end{itemize}

\section{The silver \texorpdfstring{$S$}{S}-unit coordinate and the Cayley coordinate}
\label{sec:sunit}

From here on $p=2J^2-1$ unless a divisor is explicitly
allowed, and $\Z[\sqrt2]$ is read modulo $p$ through
\begin{equation}
 \sqrt2\longmapsto2J,\qquad\text{since }(2J)^2=2p+2\equiv2\pmod p .
 \label{eq:sqrt2-reduction}
\end{equation}
This is a ring homomorphism $\Z[\sqrt2]\to\Z/p$.  It sends
$\delta:=1+\sqrt2$ to $2J+1$, $\delta^{-1}=\sqrt2-1$ to $2J-1$,
$\delta^2=3+2\sqrt2$ to $4J+3=\varepsilon$ (the silver unit of
\eqref{eq:silver-unit}), $\delta^{-2}=3-2\sqrt2$ to $3-4J$, and $2$ to a
unit, $p$ being odd.  Hence every \emph{silver $S$-unit}
$\pm\delta^n\sqrt2^{\,m}$ ($n,m\in\Z$) reduces to a unit modulo $p$; the
same holds modulo every divisor $p>1$ of $2J^2-1$, where
\eqref{eq:sqrt2-reduction} still holds.  For $\xi=x+y\sqrt2\in\Z[\sqrt2]$
the reduction is $\xi(J)=x+2yJ$.  Since $\eta_J$ depends on $q$ only
through the lens orbit $\{\pm q^{\pm1}\}$ up to sign
\eqref{eq:lens-symmetries}, we speak of the \emph{orbit} of an $S$-unit
$\xi$, meaning $\{\pm\xi^{\pm1}\}$, and write $(n,m)$ for the orbit of
$\delta^n\sqrt2^{\,m}$; every orbit has a representative in $\Z[\sqrt2]$
(replace $\xi$ by $\xi^{-1}$ if necessary).

\begin{lemma}[The four known families as $S$-unit orbits]
\label{lem:sunit-four}
Let $J\geq2$ and $p=2J^2-1$.  Modulo $p$,
\[
 J\equiv\sqrt2^{\,-1},\qquad J+1\equiv\delta/\sqrt2,\qquad
 2J-1\equiv\delta^{-1},\qquad J(J-1)\equiv-(2\delta)^{-1}.
\]
Thus the four zero families $q=J$, $J+1$, $2J-1$, $J(J-1)$ of
\cite[Thm.~7.1]{MahmoodP1Collection2026} are the orbits $(0,1)$,
$(1,-1)$, $(1,0)$, $(1,2)$ of $\sqrt2$, $\delta/\sqrt2$, $\delta$ and
$2\delta$.  The same congruences hold modulo every divisor $p>1$ of
$2J^2-1$.
\end{lemma}

\begin{proof}
$2J\cdot J=2J^2\equiv1$, so $J\equiv(2J)^{-1}\equiv\sqrt2^{\,-1}$; then
$J+1\equiv\sqrt2^{\,-1}+1=(1+\sqrt2)/\sqrt2$; $2J-1\equiv\sqrt2-1=\delta^{-1}$;
and $(4J+2)\,J(J-1)=2J(2J^2-1)-2J^2\equiv-1$, so
$J(J-1)\equiv-(4J+2)^{-1}=-(2+2\sqrt2)^{-1}$.  Only $2J^2\equiv1$ was
used.
\end{proof}

\begin{definition}[Cayley coordinate]
\label{def:cayley}
Let $p$ be odd and
\[
 D_p=\{q\in(\Z/p)^\times:\ \gcd(1-q,p)=\gcd(1+q,p)=1\}
 =\{q\in(\Z/p)^\times:\ \gcd(q^2-1,p)=1\}.
\]
For $q\in D_p$ put
\begin{equation}
 \rho=\rho(q)=\frac{1+q}{1-q}\in(\Z/p)^\times ,
 \label{eq:cayley}
\end{equation}
the ratio $A/B$ of the two coefficients in \eqref{eq:AB}, i.e.\ the
ratio of the left and right Clifford rotation angles of the generator of
$L(p,q)$.
\end{definition}

At a prime $p$, $D_p$ is the set of units $q\not\equiv\pm1$.  At a
composite $p$ it is smaller, and outside it $\rho$ is not a unit: for
$J=5$, $p=49$, the unit $q=8$ has $1-q=-7$ not invertible, and $q=6$ gives
$1+q=7$ and $\rho=28$.  On
$D_p$ the map $q\mapsto\rho$ is a M\"obius transformation, hence
injective, with inverse $q=(\rho-1)/(\rho+1)$, and $\rho\not\equiv\pm1$.
The set $D_p$ is closed under $q\mapsto q^{-1}$ (as
$1-q^{-1}=-q^{-1}(1-q)$) and $q\mapsto-q$ (which exchanges $1\pm q$);
$q\mapsto q^{-1}$ sends $\rho\mapsto-\rho$ and $q\mapsto-q$ sends
$\rho\mapsto\rho^{-1}$, so the lens orbit $\{\pm q^{\pm1}\}$ corresponds
to $\{\pm\rho^{\pm1}\}$, and orientation reversal ($\eta_J\mapsto-\eta_J$)
is $\rho\mapsto\rho^{-1}$.

\begin{proposition}[Dictionary at silver ranks]
\label{prop:cayley-dictionary}
Let $p>1$ divide $2J^2-1$ and read $\sqrt2\equiv2J$, $\delta\equiv2J+1$
as above.  The orbits below lie in $D_p$ (the mirror orbit whenever it
exists), and in the coordinate $\rho$:
\begin{enumerate}[label=(\alph*)]
\item the balanced family (b) of Theorem~\ref{thm:classification},
$q\equiv\pm(2J\pm1)$, is $\rho\in\{\pm\delta^{\pm1}\}$;
\item the orbit of Theorem~\ref{thm:C3}, $q\equiv\pm J^{\pm1}$, is
$\rho\in\{\pm\delta^{\pm2}\}$;
\item the mirror orbit $q^2\equiv-1$ is $\rho^2\equiv-1$;
\item the orbit of $q\equiv3-4J\equiv\delta^{-2}$ is
$\rho\in\{\pm\sqrt2^{\,\pm1}\}$, with $\rho(3-4J)=\sqrt2$ exactly.
\end{enumerate}
\end{proposition}

\begin{proof}
Only $2J^2\equiv1$ is used.  \emph{Membership in $D_p$.}  The residues
$2J$, $J\pm1$ and $2J\pm1$ are units: a prime dividing $J\pm1$ and
$2J^2-1$ divides $2J^2-1-2(J\mp1)(J\pm1)=1$, and
$(2J+1)(2J-1)=2(2J^2-1)+1$.  For $q=\pm(2J\pm1)$ the numbers $1\pm q$
are, up to sign, among $2J$, $2(J+1)$, $2(J-1)$; for $q=\pm J$, $\pm2J$
they are $1\pm J$, $1\pm2J$; for $q=3-4J$ they are $2(2J-1)$ and
$-4(J-1)$, and for its inverse $3+4J$ they are $-2(2J+1)$ and $4(J+1)$;
and for $q^2\equiv-1$, $q^2-1\equiv-2$.  All are units, $p$ being odd,
and $D_p$ is closed under $q\mapsto\pm q^{\pm1}$.  (A check over every
divisor $p>1$ of $2J^2-1$, $2\leq J\leq100$, found no exception
\cite{LCrecord}.)
(a) For $q=2J+1$:
$\rho=(2J+2)/(-2J)=-(J+1)J^{-1}=-(J+1)\cdot2J=-(2J^2+2J)\equiv-(2J+1)=-\delta$.
(b) For $q=J$: $(3+4J)(1-J)=3+J-4J^2\equiv1+J$, so
$\rho=(1+J)/(1-J)=3+4J=\delta^2$.
(c) $(1+q)^2+(1-q)^2=2(1+q^2)$, so $\rho^2\equiv-1$ iff $q^2\equiv-1$.
(d) For $q=3-4J$: $\rho=(4-4J)/(4J-2)=(2-2J)/(2J-1)$, and
$2J(2J-1)=4J^2-2J\equiv2-2J$, so $\rho=2J=\sqrt2$.  The orbit statements
follow from the transformation rules above.
\end{proof}

\begin{observation}[Silver primes in the coordinate $\rho$]
\label{obs:rho-statistics}
The following was found in exploratory computations, every prime $p\leq673$ that has
a silver rank being taken at its minimal silver rank $J$ (the smaller
root of $2J^2\equiv1$).  ``Sporadic'' means: not in the three orbits of
(a).
\begin{enumerate}[label=(\alph*)]
\item \textup{(Universal zero orbits.)}  The orbits that are zeros at
every silver prime examined are the two $S$-unit orbits of
Proposition~\ref{prop:cayley-dictionary}(a),(b), namely family (b),
witnessed by level balance, and the orbit of Theorem~\ref{thm:C3},
witnessed by the inert class $R^2L^2$ of
\cite[Thm.~8.2(viii)]{LensChirality2026}; and the mirror orbit, witnessed
by orientation reversal, which is not an $S$-unit orbit and exists
exactly at the silver primes $p\equiv1\pmod8$.  Among the $S$-units
$q=s\cdot2^c\sqrt2^{\,a}\delta^{b}$ with $s=\pm1$, $|c|\leq2$, $|a|\leq3$,
$|b|\leq4$, exactly the $28$ expressions lying in those two orbits are
zeros at all $57$ silver primes $17\leq p\leq673$; among
$\rho=\pm\delta^k$, $1\leq k\leq14$, exactly $k=1,2$ are (the next best,
$k=14$, at $16$ of $56$ primes where it is defined).  These searches were
first run during exploration on smaller prime ranges.  With the expression
ranges recovered from the retained scripts, repetition over the stated
prime range tested and confirmed the same predicted universal-zero sets
\cite{LCrecord}.
\item \textup{(Sporadic zeros, primes $<400$.)}  The sporadic zeros show
no structure in $\rho$ on the unit group or in lattice size: on the
subgroup $\langle-1,\delta\rangle$ they occur at the subgroup's density
(for instance $0$ of $3$ at $p=41$ and $0$ of $12$ at $p=199$); their
rate is flat, $16$--$20\%$, across height buckets of the
$\Z[\sqrt2]$-lattice; and the hypothesis that the cofactor norm lies in
$\{m,2m\}$ is refuted (rates $5/51$ and $8/45$ against $229/1384$ for
all other classes).
\item \textup{(One signal.)}  Among sporadic cases, lifts of norm $2$
($\rho\equiv\sqrt2\cdot$unit) are zeros at $40/137=29\%$ against $15\%$
otherwise (primes $\leq400$); the excess sits at $\rho=\pm\sqrt2$
exactly, i.e.\ at $q\equiv3-4J\equiv2(J-1)^2$.  At $\rho=\sqrt2$ the
zero occurs at $0$ of the $28$ primes $p\equiv1\pmod8$ and at $18$ of the
$29$ primes $p\equiv7\pmod8$ with $p\leq673$; and at the primes
$p=2J^2-1$ with $3\leq J\leq18$ every row fits
$\eta_J(L(p,3-4J))=J-1$ for odd $J$ and $0$ for even $J$.
\end{enumerate}
\end{observation}

\begin{remark}[Why the mechanism of {\cite[Sec.~7]{LensChirality2026}} does not extend]
\label{rem:kill}
The proof of Theorem~\ref{thm:C3} in \cite[Sec.~7]{LensChirality2026} succeeded because at $q=J$ the summand
of the trace form \cite[eq.~(19)]{LensChirality2026} collapsed \cite[Lems.~7.2--7.3]{LensChirality2026} to a single
two-factor Fourier--Dedekind sum with the quarter shift of
\cite[Lem.~7.4]{LensChirality2026}, and that shift produced one untwisted
Rademacher symbol.  For general $q$ this does not happen.  In the
variable $\zeta=\e^{2\pi ik/p}$ the denominators of the two terms of
the trace form are, up to constants and roots of unity,
$(\zeta^{a}-1)(\zeta^{q-1}-1)(\zeta^{q+1}-1)$ with $a\in\{1,q\}$, so
$p\,\eta_J$ is a combination of three-factor Fourier--Dedekind sums.
With $A=\zeta^{q-1}-1$, $B=\zeta^{q+1}-1$ and $C=\zeta^2-1$ one has
$B-A=\zeta^{q-1}C$, hence the cocycle identity
\[
 \frac1{AB}=\frac1C\Bigl(\frac1A-\frac1B\Bigr)-\frac1B ,
\]
which reduces them to two-factor sums; but the resulting shifts are $Jq$
and $(J+1)q$, not quarter shifts, and the sums are twisted Dedekind sums
(Rademacher symbols at congruence level) rather than one untwisted
$\Psi$.  Thus the conjugacy mechanism does not by itself close the
classification at silver ranks.  Nothing about the twisted sums is
claimed here.
\end{remark}

\begin{conjecture}[Nonvanishing at $q=3-4J$ for primes $p\equiv1\pmod8$; refuted]
\label{conj:C5}
For every prime $p\equiv1\pmod8$ with minimal silver rank $J$,
$\eta_J(L(p,3-4J))\neq0$.
\end{conjecture}

This was formulated, together with Observation~\ref{obs:rho-statistics}(c),
before the $28$ primes $p\equiv1\pmod8$ with $761\leq p\leq1553$ were
computed.  It is false: $27$ of them give $\eta_J\neq0$, but $p=1129$,
$J=203$ gives $\eta_{203}(L(1129,\,3-4\cdot203))=0$.  Together with the
exploratory prime sample, the zero at $\rho=\sqrt2$ occurs at $1$ of $56$ primes
$p\equiv1\pmod8$ and at $18$ of $29$ primes $p\equiv7\pmod8$; the split
by $p\bmod8$ is a statistic, not a law, and no replacement conjecture is
made.  (The complementary control ``a zero at every silver prime'' fails
at $p\equiv1\pmod8$.)

\section{The fifth family at \texorpdfstring{$p=2J^2-1$}{p = 2J\textasciicircum 2 - 1}}
\label{sec:fifth}

The orbit $\rho=\sqrt2$, i.e.\ $q\equiv3-4J\equiv\delta^{-2}$, is the
one place where Observation~\ref{obs:rho-statistics} showed structure.
Its value at $p=2J^2-1$ was conjectured from the exploratory rows
$3\leq J\leq18$ with $p$ prime, confirmed on every $J\in[19,60]$ with
$p$ prime in a subsequent test of that prediction ($19$ rows,
Appendix~\ref{app:verification}), and then proved for every $J$, prime
or composite $p$, by the four-term method of
\cite[Thm.~7.1]{MahmoodP1Collection2026}.  Write the master formula
\eqref{eq:master} as
\begin{equation}
 \eta_J(L(p,q))=-4J\,s(q,p)+R_J(p,q)-2\Delta_J(p,q)-L_J(p,q),
 \label{eq:master-four-terms}
\end{equation}
with the residue term and the triangular sum
\[
 R_J=\frac2p\sum_{a=1}^{J-1}(J-a)\bigl[\rho_p(aq)+\rho_p(aq^*)\bigr],\qquad
 \Delta_J=\!\!\sum_{\substack{a,b\geq1\\a+b\leq J-1}}\!\!(J-a-b)
 \bigl[\delta_p(aq-b)-\delta_p(aq+b)\bigr],
\]
and the low-mode term $L_J$ of Section~\ref{sec:conventions}.
Throughout, $q^*$ is an integer inverse of $q$ modulo $p$,
$\ind_{\mathcal P}\in\{0,1\}$ is the indicator of a proposition
$\mathcal P$, and
\begin{equation}
 \mathcal T(x,y)=\frac1{12}\Bigl(\frac xy+\frac yx+\frac1{xy}\Bigr)-\frac14,
 \qquad s(x,y)+s(y,x)=\mathcal T(x,y)\quad(x,y\geq1,\ \gcd(x,y)=1),
 \label{eq:reciprocity}
\end{equation}
is Dedekind reciprocity.

\begin{theorem}[The fifth family]
\label{thm:C4}
Let $J\geq2$, $p=2J^2-1$ (prime or not) and $q\equiv3-4J\pmod p$, whose
least positive residue is $2(J-1)^2$ and whose inverse is
$q^*\equiv3+4J$.  Then
\begin{equation}
 \eta_J(L(p,q))=(J-1)\,\ind_{J\text{ odd}} .
 \label{eq:C4}
\end{equation}
In particular $q\equiv3-4J$ is a zero of $\eta_J$ exactly for even $J$,
and the same value holds on the orbit $\pm q^{\pm1}$ up to the sign of
$q$.
\end{theorem}

\begin{proof}
\emph{Units.}  $(3-4J)(3+4J)=9-16J^2=1-8p$, so $q^*\equiv3+4J=:h$; and
$8p=(4J+3)(4J-3)+1$ gives $\gcd(h,p)=1$.

\emph{Step 1 (the Dedekind sum).}  $s(q,p)=s(q^*,p)=s(h,p)$.  By
\eqref{eq:reciprocity}, $s(h,p)=-s(p,h)+\mathcal T(h,p)$.  Since
$8p\equiv1\pmod h$, $s(p,h)=s(8,h)$, and reciprocity again gives
$s(8,h)=-s(h,8)+\mathcal T(8,h)$.  Using $h^2+1-8p=6h$ the two
reciprocity terms collapse to $\mathcal T(h,p)-\mathcal T(8,h)=\frac1{2p}-\frac1{16}$,
so
\[
 s(h,p)=s(h,8)+\frac1{2p}-\frac1{16}.
\]
Here $h\equiv3\pmod8$ for even $J$ and $h\equiv7\pmod8$ for odd $J$, and
$s(3,8)=\frac1{16}$, $s(7,8)=-\frac7{16}$.  Hence
\begin{equation}
 s(q,p)=\frac1{2p}-\frac{\ind_{J\text{ odd}}}2 ,
 \label{eq:C4-dedekind}
\end{equation}
which is the value of Corollary~\ref{cor:S} at $\varepsilon=q^*$, since
$p\equiv7\pmod8$ for even $J$ and $p\equiv1\pmod8$ for odd $J$.

\emph{Step 2 (the residue term).}  For $1\leq a\leq J-1$ both $a(4J+3)$
and $a(4J-3)$ are less than $2p$, so each residue wraps at most once; the
wrap thresholds are
\[
 \frac p{4J+3}=\frac J2-\frac38+\frac1{8(4J+3)},\qquad
 \frac p{4J-3}=\frac J2+\frac38+\frac1{8(4J-3)} .
\]
With $aq\equiv-a(4J-3)$, $aq^*=a(4J+3)$, $\rho_p(-k)=-\rho_p(k)$ and
$\rho_p(k)=p-2\langle k\rangle_p$, an index $a$ at which neither or both
residues wrap contributes $\rho_p(aq)+\rho_p(aq^*)=-12a$; the only mixed
index is $a=J/2$ for even $J$, where $aq^*$ wraps and $a(4J-3)$ does
not, and it contributes $2p-12a$.  For odd $J$ there is no mixed index.
Therefore
\begin{equation}
 R_J=-\frac{4J(J^2-1)}p+2J\,\ind_{J\text{ even}} .
 \label{eq:C4-residue}
\end{equation}

\emph{Step 3 (the triangular sum vanishes).}  A term $aq\pm b\equiv0$
with $a,b\geq1$, $a+b\leq J-1$ means $a(4J-3)\mp b\in\{p,2p\}$.  Each
case confines $a$ to an interval of length less than $\frac12$ containing
no integer, or forces $a>J-1$: $a(4J-3)+b=p$ needs
$\frac J2+\frac J{4(J-1)}\leq a<\frac J2+\frac38$; $a(4J-3)-b=p$ needs
$\frac J2+\frac38<a\leq\frac{J+1}2-\frac1{2(2J-1)}$; and the value $2p$
needs $a>J-1$.  Hence $\Delta_J=0$ (for $J=2$ the sum is empty).

\emph{Step 4 (the low-mode term).}  As in Section~\ref{sec:conventions},
a weight pair $(w_L,w_R)\in W_{a-1}\times W_{b-1}$ is invariant iff
$(1+\tilde q)w_L+(1-\tilde q)w_R\equiv0\pmod{2p}$; with
$u=(w_L-w_R)/2$, $v=(w_L+w_R)/2\in\Z$ (as $w_L\equiv w_R\pmod2$) this
reads $v+qu\equiv0\pmod p$, i.e.\ $v\equiv(4J-3)u$, and in the low-mode
range $|u|,|v|\leq J-1$.  Writing $v=(4J-3)u-\ell p$: $\ell=0$ forces
$u=v=0$; $\ell=1$ needs $(4J-3)u\in[p-J+1,p+J-1]$, i.e.\
$u\in\bigl[\frac J2+\frac{J/2}{4J-3},\ \frac J2+\frac{5J/2-2}{4J-3}\bigr]$,
which contains an integer only for odd $J$, namely $u=\frac{J+1}2$,
$v=\frac{J-1}2$, i.e.\ $(w_L,w_R)=(J,-1)$; $\ell=2$ would need $u>J-1$;
negative $\ell$ give the negatives of these points.  The solution
$(0,0)$ lies in $W_{a-1}\times W_{b-1}$ iff $a$ and $b$ are both odd, and
$J\pm m$ have the same parity, so it contributes equally to
$I_{J-m,J+m}$ and $I_{J+m,J-m}$ and cancels in $L_J$.  The solutions
$\pm(J,-1)$ never lie in $W_{J-m-1}\times W_{J+m-1}$ (as
$|w_L|=J>J-m-1$), and they lie in $W_{J+m-1}\times W_{J-m-1}$ exactly
when $m$ is odd and $1\leq m\leq J-2$ (the parity and size conditions of
the two weight sets), which requires $J$ odd.  There are $\frac{J-1}2$
such $m$, each contributing $2$ to $I_{J+m,J-m}$, so
\begin{equation}
 L_J=-(J-1)\,\ind_{J\text{ odd}} .
 \label{eq:C4-lowmode}
\end{equation}

\emph{Assembly.}  By \eqref{eq:master-four-terms} and
\eqref{eq:C4-dedekind}--\eqref{eq:C4-lowmode},
\begin{align*}
 \eta_J&=-4J\Bigl(\frac1{2p}-\frac{\ind_{J\text{ odd}}}2\Bigr)
 -\frac{4J(J^2-1)}p+2J\,\ind_{J\text{ even}}+(J-1)\,\ind_{J\text{ odd}}\\
 &=-\frac{2J}p\bigl(1+2(J^2-1)\bigr)+2J+(J-1)\,\ind_{J\text{ odd}},
\end{align*}
and $1+2(J^2-1)=p$, which gives \eqref{eq:C4}.  The orbit statement is
\eqref{eq:lens-symmetries}.
\end{proof}

\begin{remark}
\label{rem:C4-history}
(i) Formula \eqref{eq:C4} was first conjectured from the exploratory rows
$3\leq J\leq18$ with $p$ prime (for instance $\eta_J(L(p,3-4J))=2,0,6,0$
at $p=17,31,97,127$) and predicted the value on the $19$ rows
$J\in\{21,22,24,25,28,34,36,38,39,41,42,43,45,46,49,50,52,56,59\}$
($p$ prime) before they were computed: the odd $J$ gave
$20,24,38,40,42,44,48,58$ and every even $J$ gave $0$.  The proof adds
the composite levels, on which nothing had been computed; for instance
$J=5$, $p=49$, $\eta_5=4$; $J=9$, $p=161$, $\eta_9=8$; $J=12$, $p=287$,
$\eta_{12}=0$.
(ii) On every row of the prime-level verification ($J\leq60$) with
$J\geq6$, the orbit of $3-4J$ is distinct from the four orbits of
Lemma~\ref{lem:sunit-four}, from family (b), from the orbit of
Theorem~\ref{thm:C3} and from the mirror orbit; at $J=4$ it coincides
with the orbit of $J(J-1)$ and at $J=2$ with every family.  On the rows
examined it is therefore a fifth family at $p=2J^2-1$ for $J\geq6$; a
proof of distinctness for all $J$ is not written out here.
(iii) A control for the verification of Theorem~\ref{thm:C4}, which
assumed that the neighbouring residue $q\equiv4-4J$ is \emph{not} a zero
on some even-$J$ row, failed: $\eta_J(L(2J^2-1,4-4J))=0$ on all eight
even $J\in\{12,14,16,20,26,30,32,40\}$ that were checked (all with
composite $p$).  In $S$-unit terms $4-4J\equiv2\sqrt2\,\delta^{-1}$, the
orbit $(1,-3)$; its behaviour is covered only by the period law of
Section~\ref{sec:period-law}, and it is recorded as an observation, not
a result.
\end{remark}

\section{The Pell ladder}
\label{sec:pell}

Theorem~\ref{thm:C4} concerns the orbit of $\delta^2$; the third family
of \cite[Thm.~7.1]{MahmoodP1Collection2026} is the orbit of $\delta$
(Lemma~\ref{lem:sunit-four}).  Both are rungs of one ladder.  Write
\begin{equation}
 \begin{gathered}
 \delta^n=(1+\sqrt2)^n=a_n+P_n\sqrt2,\\
 (a_n)_{n\geq0}=1,1,3,7,17,41,99,\ldots,\qquad
 (P_n)_{n\geq0}=0,1,2,5,12,29,70,\ldots,
 \end{gathered}
 \label{eq:pell-numbers}
\end{equation}
so that $a_{n+1}=2a_n+a_{n-1}$, $P_{n+1}=2P_n+P_{n-1}$ and
$a_n=P_n+P_{n-1}$, and put $\nu_n=(-1)^n$, the norm of $\delta^n$.  For
$J\geq2$ and $p=2J^2-1$ let
\begin{equation}
 q_n:=a_n+2P_nJ ,
 \label{eq:qn}
\end{equation}
the reduction of $\delta^n$ by \eqref{eq:sqrt2-reduction}, taken as a
positive integer.  Thus $q_1=2J+1$ (family (b), the third family of
\cite[Thm.~7.1]{MahmoodP1Collection2026}) and
$q_2=4J+3\equiv(3-4J)^{-1}$ (Theorem~\ref{thm:C4}).  Finally
$M_n:=2P_n^2$.

\begin{theorem}[The Pell ladder]
\label{thm:P}
Let $J\geq2$, $p=2J^2-1$ and $q_n$ as in \eqref{eq:qn}.
\begin{enumerate}[label=(\roman*)]
\item If $n\geq1$ and $P_n\mid J$, then
\[
 \eta_J(L(p,q_n))=\begin{cases}0,&n\text{ odd},\\[2pt]
 \tfrac23\,P_{n-2}\,J,&n\text{ even}.\end{cases}
\]
\item If $n\geq2$ and $J\equiv1\pmod{P_n}$, then
\[
 \eta_J(L(p,q_n))=\gamma_n\,\frac{J-1}{P_n},\qquad
 \gamma_n=\begin{cases}2P_{n-1}^2,&n\text{ even},\\[2pt]
 \tfrac43\,P_{n-1}a_{n-1},&n\text{ odd}.\end{cases}
\]
\end{enumerate}
The same values hold on the orbit $\pm q_n^{\pm1}$, with the sign of $q$
\eqref{eq:lens-symmetries}.
\end{theorem}

\begin{corollary}
\label{cor:P}
\begin{enumerate}[label=(\alph*)]
\item On the multiples of $P_n$, the orbit of $\delta^n$ is a zero of
$\eta_J$ iff $n$ is odd or $n=2$.
\item For positive $n$, the orbit of $\delta^n$ vanishes at every rank
$J\geq2$, or even at all sufficiently large ranks, iff $n=1$ (the third family of
\cite[Thm.~7.1]{MahmoodP1Collection2026}); for $n\geq2$ part (ii) gives
$\eta_J>0$ on the infinite class $J\equiv1\pmod{P_n}$.
\item $n=2$ recovers Theorem~\ref{thm:C4} on both parities: $0$ for even
$J$ by (i), and $\gamma_2(J-1)/2=J-1$ for odd $J$ by (ii).
\end{enumerate}
\end{corollary}

\begin{proof}
(a) For odd $n$ the value in (i) is $0$; for even $n$ it is $0$ iff
$P_{n-2}=0$, i.e.\ $n=2$, since $P_{n-2}\geq2$ for even $n\geq4$.  (b) For
$n=1$, $P_1=1$ divides every $J$ and (i) gives $0$; for $n\geq2$,
$\gamma_n>0$.  (c) $P_2=2$ and $\gamma_2=2P_1^2=2$.
\end{proof}

\begin{lemma}[The Dedekind term at every rank]
\label{lem:S}
For $n\geq1$, $J\geq1$, $p=2J^2-1$ and $q_n$ as in \eqref{eq:qn},
\begin{equation}
 s(q_n,p)=\mathcal T(q_n,p)-\nu_n\,\mathcal T(M_n,q_n)
 +\nu_n\,s\bigl(q_n\bmod M_n,\,M_n\bigr),
 \label{eq:lemma-S}
\end{equation}
and $q_n\bmod M_n=a_n+2P_n\,(J\bmod P_n)$ depends on $J$ only modulo
$P_n$.
\end{lemma}

\paragraph{Pell facts.}
Fix $n\geq1$ and write $a=a_n$, $b=P_n$, $\nu=\nu_n$, $M=M_n=2b^2$.
From $\delta^{n\pm1}=\delta^n\delta^{\pm1}$, $\delta^{-2}=3-2\sqrt2$ and
the norm:
\begin{enumerate}[label=(F\arabic*)]
\item $a^2-M=\nu$.
\item $a_{n-1}=2b-a\geq1$, $P_{n-1}=a-b\geq0$, $a_{n+1}=a+2b$,
$P_{n-2}=3b-2a$ for $n\geq2$, and $P_{n-3}=5a-7b$ for $n\geq3$.
\item $a_{n+1}a_{n-1}=M-\nu$.
\item $b\leq a<2b$; $a$ is odd and $\gcd(a,2b)=1$; and $2a\leq3b$ for
every $n\geq1$, with equality only at $n=2$ (for $n\geq2$ this is
$P_{n-2}\geq0$; at $n=1$ it follows from $a=b=1$).
\item With $t:=2bJ-a$,
\begin{equation}
 M\,p=q_n\,t+\nu ,
 \label{eq:F5}
\end{equation}
because $M(2J^2-1)=(2bJ)^2-a^2+(a^2-2b^2)$.  Hence
$q_n\cdot(-\nu t)\equiv1\pmod p$, so $q_n^*\equiv-\nu t$ and
$\rho_p(\alpha q_n^*)=-\nu\,\rho_p(\alpha t)$ for every integer
$\alpha$.
\end{enumerate}
(F2) and (F3) are read off from $(a+b\sqrt2)(\sqrt2\mp1)$,
$(a+b\sqrt2)(3-2\sqrt2)$ and $(a+2b)(2b-a)=4b^2-a^2$; (F4) follows from
(F1) and (F2) and the recurrence ($a_n$ is odd).

\begin{proof}[Proof of Lemma~\ref{lem:S}]
Write $q=q_n$.  By \eqref{eq:F5}, $Mp\equiv\nu\pmod q$, so
$p\equiv\nu M^{-1}\pmod q$.  Moreover $\gcd(M,q)=1$: $q=a+2bJ$ is odd,
and modulo each prime factor of $b$ one has $q\equiv a\not\equiv0$ by
(F4).  Since $s(h^{-1},k)=s(h,k)$ and $s(-h,k)=-s(h,k)$, this gives
$s(p,q)=\nu\,s(M,q)$.  Two applications of \eqref{eq:reciprocity},
$s(q,p)=\mathcal T(q,p)-s(p,q)$ and $s(M,q)=\mathcal T(M,q)-s(q,M)$, give
\eqref{eq:lemma-S}, as $s(q,M)$ depends only on $q\bmod M$.  Finally
$2bJ\equiv2b(J\bmod b)\pmod{2b^2}$ and $a+2b(b-1)<2b^2$ by (F4).
\end{proof}

\paragraph{What Lemma~\ref{lem:S} does and does not give.}
With $r:=a+2b\,(J\bmod b)$ the identity \eqref{eq:lemma-S} simplifies
exactly to
\begin{equation}
 s(q_n,p)=\frac a{6p}-\frac a{6M}+s(r,M)\ \ (n\text{ even}),\qquad
 s(q_n,p)=\frac{bJ}{3p}+\frac J{6b}-\frac12-s(r,M)\ \ (n\text{ odd}).
 \label{eq:lemma-S-explicit}
\end{equation}
Only the lookup $s(r,M)$ depends on $J$ through $J\bmod P_n$; the
rational terms do not vanish, and they are not eventually
quasi-polynomial in $J$ (for $n=1$ there is a single residue class, yet
$s(5,7)=-\tfrac1{14}$ at $J=2$ and $s(7,17)=\tfrac1{17}$ at $J=3$).
Lemma~\ref{lem:S} is thus an exact reduction of the Dedekind term to a
rational function of $J$ plus a finite periodic lookup; it is not a
proof of a ``Dedekind-term part'' of Conjecture~\ref{conj:PL} below.
The cancellation of the rational remainders happens only in the
complete $\eta_J$: on the two slices of Theorem~\ref{thm:P} the residue
term $R_J$ absorbs them, and in the sectors of
Section~\ref{sec:period-law-proved} the cancellation is proved through
the integrality of $\eta_J$ (Theorem~\ref{thm:PL}).

\begin{proof}[Proof of Theorem~\ref{thm:P}]
We evaluate the four terms of \eqref{eq:master-four-terms} at $q=q_n$,
using \eqref{eq:F5} to control every residue; $\alpha$ denotes the
summation index of $R_J$ and $(\alpha,\beta)$ that of $\Delta_J$.

\medskip\noindent\emph{Slice (i): $J=bk$, $k\geq1$, $J\geq2$.}  Here
$q=Mk+a$, $t=Mk-a$ and $p=Mk^2-1$.

\emph{Step 1 (the Dedekind sum).}  $q\bmod M=a$, and by (F1)
$M\equiv-\nu\pmod a$, so $s(a,M)=\mathcal T(a,M)+\nu\,s(1,a)$ with
$s(1,a)=\frac{(a-1)(a-2)}{12a}$.  Lemma~\ref{lem:S} then simplifies to
\begin{equation}
 s(q,p)=\frac{Mk+(k-3)p}{6p}\ \ (n\text{ odd}),\qquad
 s(q,p)=\frac a{6p}+\frac{a-3}6\ \ (n\text{ even}).
 \label{eq:P-i-dedekind}
\end{equation}
At $n=2$ this is $\frac1{2p}$, the even-$J$ value in
\eqref{eq:C4-dedekind}.

\emph{Step 2 (residues).}  For $1\leq\alpha\leq J-1$ write $\alpha=ck+r$
with $0\leq c\leq b-1$ and $0\leq r\leq k-1$.  Then
\[
 \alpha q=cp+c(ak+1)+rq,\qquad \alpha t=cp-c(ak-1)+rt .
\]
The residue $\langle\alpha q\rangle_p=c(ak+1)+rq$ never wraps: its
largest value is $p-b\,a_{n-1}k+(b-a)<p$.  The residue
$\langle\alpha t\rangle_p$ equals $rt-c(ak-1)\in(0,p)$ for $r\geq1$, and
$p-c(ak-1)$ (one wrap) for $r=0$, $c\geq1$.  Hence
\[
 \rho_p(\alpha q)=p-2c(ak+1)-2rq,\qquad
 \rho_p(\alpha t)=p-2rt+2c(ak-1)-2p\,\ind_{r=0},
\]
and summing $(J-\alpha)[\rho_p(\alpha q)-\nu\rho_p(\alpha t)]$ over
$\alpha$ (Faulhaber sums in $c$ and $r$) gives
\begin{equation}
 R_J=4J\,s(q,p)\ \ (n\text{ odd}),\qquad
 R_J=4J\,s(q,p)+\tfrac23P_{n-2}J\ \ (n\text{ even}).
 \label{eq:P-i-residue}
\end{equation}
For odd $n$ this is $R_J=4Js$, the cancellation pattern of the first
three families of \cite[Thm.~7.1]{MahmoodP1Collection2026}; for even $n$
it reads $R_J=2b(b-1)k-\frac{4a}{3p}(J^3-J)$.

\emph{Step 3 (the triangular sum is empty).}  A term $\alpha q\equiv\beta$
with $1\leq\beta\leq J-1-\alpha$ would need
$\langle\alpha q\rangle_p\leq J-1$; but $\langle\alpha q\rangle_p\geq q>J$
when $r\geq1$ and $\langle\alpha q\rangle_p\geq ak+1>bk=J$ when $r=0$.  A
term $\alpha q\equiv-\beta$ would need $p-\langle\alpha q\rangle_p\leq J-1$;
but $p-\langle\alpha q\rangle_p\geq b\,a_{n-1}k+(a-b)\geq J$.  So
$\Delta_J=0$.

\emph{Step 4 (the low-mode term vanishes).}  In the coordinates $u,v$ of
the proof of Theorem~\ref{thm:C4}, the invariant low-mode weights are the
lattice points with $|u|,|v|\leq J-1$ and $v+qu\equiv0\pmod p$.  For
$1\leq u\leq J-1$ such a point needs $\langle uq\rangle_p\leq J-1$ or
$\langle uq\rangle_p\geq p-J+1$, both excluded by Step~3; negative $u$
follow by symmetry, and $u=0$ forces $v=0$.  Only $(0,0)$ remains, and it
cancels in $L_J$ as before.  So $L_J=0$.

\emph{Assembly.}  $\eta_J=-4Js+R_J$, which is $0$ for odd $n$ and
$\frac23P_{n-2}J$ for even $n$; this is (i).

\medskip\noindent\emph{Slice (ii): $n\geq2$, $J=bk+1$, $k\geq1$.}  Write
$a_-=a_{n-1}$ and $a_+=a_{n+1}$, so that $M=a_+a_-+\nu$ by (F3).  Then
$q=Mk+a_+$, $t=Mk+a_-$, $p=Mk^2+4bk+1$, and with
\[
 g:=a_-k+1,\qquad h:=a_+k+1,\qquad p=kq+g=kt+h .
\]

\emph{Step 1 (the Dedekind sum).}  $q\bmod M=a_+$ (as $a_+<M$ for
$b\geq2$), and $M\equiv\nu\pmod{a_+}$ by (F3), so
$s(a_+,M)=\mathcal T(a_+,M)-\nu\,s(1,a_+)$.  Lemma~\ref{lem:S} then
simplifies to
\begin{equation}
 s(q,p)=-\frac{a(J^2-1)}{3p}\ \ (n\text{ even}),\qquad
 s(q,p)=\frac{J-1}{6b}-\frac{b(J-1)(2J+1)}{3p}\ \ (n\text{ odd}).
 \label{eq:P-ii-dedekind}
\end{equation}
At $n=2$ this is $\frac1{2p}-\frac12$, the odd-$J$ value in
\eqref{eq:C4-dedekind}.

\emph{Step 2 (residues).}  For $1\leq\alpha\leq J-1=bk$ write
$\alpha=ck+r$ with $0\leq c\leq b$ and $0\leq r\leq k-1$, where $c=b$
occurs only with $r=0$.  Then $\alpha q=cp+rq-cg$ and
$\alpha t=cp+rt-ch$.  The residue $\langle\alpha q\rangle_p$ equals
$rq-cg\in(0,p)$ for $r\geq1$ and $p-cg$ for $r=0$.  For
$\langle\alpha t\rangle_p$ put $c^*:=a_-+\ind_{n\text{ even}}$; then
$\langle\alpha t\rangle_p=rt-ch$ for $r\geq2$ and for $r=1$ with $c<c^*$;
$=p+rt-ch$ for $r=1$ with $c\geq c^*$; and $=p-ch$ for $r=0$.  The sign
at $r=1$ comes from
\[
 t-ch=\bigl(a_+(a_--c)+\nu\bigr)k+(a_--c) .
\]
For $k=1$ there is no row $r=1$; instead $\langle\alpha t\rangle_p=2p-ch$
exactly for $c^*+1\leq c\leq b$, because $c^*h\leq p<(c^*+1)h$ at $k=1$.
That correction equals the $r=1$ correction of the formula for $k\geq2$,
so one closed form covers every $k\geq1$.  Summing as in slice (i),
\begin{equation}
 R_J=4J\,s(q,p)+\kappa_nk,\qquad
 \kappa_n=\begin{cases}
 2P_{n-1}(P_{n-1}-1),&n\text{ even},\\
 2P_{n-1}(P_{n-1}-1)-\tfrac23P_{n-1}P_{n-3},&n\text{ odd }(n\geq3).
 \end{cases}
 \label{eq:P-ii-residue}
\end{equation}

\emph{Step 3 (the triangular sum).}  A term $\alpha q\equiv+\beta$ is
impossible: $\langle\alpha q\rangle_p>J$ when $r\geq1$ and
$\langle\alpha q\rangle_p\geq p-bg>J$ when $r=0$.  For
$\alpha q\equiv-\beta$ with $r\geq1$, $p-\langle\alpha q\rangle_p\geq g+q>J$.
For $r=0$, $\alpha=ck$ and $\beta=cg$, and $\alpha+\beta\leq J-1$ reads
$c(1+a_-)k+c\leq bk$; this holds only for $c=1$, since $2(1+a_-)>b$ by
(F4), and it needs $P_{n-1}=b-a_-\geq2$.  So there is exactly one pair,
$(\alpha,\beta)=(k,\,a_-k+1)$, of type $\alpha q+\beta\equiv0$, with
weight $J-\alpha-\beta=(P_{n-1}-1)k$, and
\begin{equation}
 \Delta_J=-(P_{n-1}-1)\,k ,
 \label{eq:P-ii-delta}
\end{equation}
which is $0$ for $n=2$.

\emph{Step 4 (the low-mode term).}  By the same bounds the only nonzero
lattice points with $|u|,|v|\leq J-1$ and $v+qu\equiv0$ are
$\pm(k,\,a_-k+1)$; the point with $c=2$ would need $2(a_-k+1)\leq bk$,
false by (F4).  In weights,
$(w_L,w_R)=\bigl((1+a_-)k+1,\ (a_--1)k+1\bigr)$, both odd since $a_-$ is
odd.  This point lies in $W_{J-m-1}\times W_{J+m-1}$ for
$m\equiv J\pmod2$ with $1\leq m\leq(P_{n-1}-1)k-1$, and in
$W_{J+m-1}\times W_{J-m-1}$ for $m\equiv J\pmod2$ with
$1\leq m\leq(P_{n-1}+1)k-1$.  The two ranges differ by exactly $k$ values
of $m$ (an interval of $2k$ consecutive integers, or $[1,2k-1]$ for
$n=2$ with $J$ odd).  Counting both points $\pm$ with the sign of $L_J$,
\begin{equation}
 L_J=-2k .
 \label{eq:P-ii-lowmode}
\end{equation}

\emph{Assembly.}  By \eqref{eq:master-four-terms} and
\eqref{eq:P-ii-residue}--\eqref{eq:P-ii-lowmode},
\[
 \eta_J=\bigl(\kappa_n+2(P_{n-1}-1)+2\bigr)k=\gamma_nk ,
\]
using $P_{n-3}=P_{n-1}-2P_{n-2}$ and $a_{n-1}=P_{n-1}+P_{n-2}$ for odd
$n$; and $k=(J-1)/P_n$.  This is (ii).  The orbit statement is
\eqref{eq:lens-symmetries}.
\end{proof}

\begin{remark}[Certificate for the finite sums]
\label{rem:faulhaber-certificate}
The closed forms \eqref{eq:P-i-dedekind}, \eqref{eq:P-i-residue},
\eqref{eq:P-ii-dedekind} and \eqref{eq:P-ii-residue} are identities
between rational functions of $(a,b,k)$ modulo the relation
$a^2=2b^2+\nu$; they follow from the residue structure of Steps~1--2 by
Faulhaber sums, which are not written out above.  They were checked by
two independent routes: with sympy over exact rationals (polynomial
remainder modulo $a^2-2b^2-\nu$, no numerical simplification) and with
the local Wolfram Engine~14.3 (\texttt{PolynomialRemainder}), four
identity sets on each route, all true.  Every step was also checked
numerically on every row of Appendix~\ref{app:verification}.  The proof
of Theorem~\ref{thm:P} is certificate-assisted at exactly this point;
the residue analysis, the emptiness and one-pair statements of Step~3,
and the lattice counts of Step~4 are as written.  (The certificates of
Section~\ref{sec:period-law-proved} also have both routes; see
Remark~\ref{rem:period-certificates}.)
\end{remark}

\begin{remark}[Verification, and the size of the rungs]
\label{rem:P-verification}
Theorem~\ref{thm:P} was verified term by term on $212$ rows, $81$ of
which were never computed before the statement was fixed, up to $n=16$;
the largest are $J=2P_{16}=941{,}664$, $p=1{,}773{,}462{,}177{,}791$ in
slice (i), with $\eta_J=50{,}713{,}000{,}832=\frac23P_{14}J$, and
$J=941{,}665$ in slice (ii), with $\Delta_J=-390{,}048=-(P_{15}-1)\cdot2$,
$L_J=-4$ and $\eta_J=152{,}139{,}002{,}500=2\gamma_{16}$.
Lemma~\ref{lem:S} reproduces the value of Corollary~\ref{cor:S} on every
$n=2$ row.  A mutated modulus, ``$\eta_J=0$ on $J\equiv0\pmod{a_n}$'',
fails at every $n$ tested, for instance $\eta_J=284,-1144,1152$ at
$J=707,714,721$ ($n=3$).
\end{remark}

\section{The period law}
\label{sec:period-law}

Theorem~\ref{thm:P} says that on the orbit of $\delta^n$ the value of
$\eta_J$ is governed by $J$ modulo $P_n$, the coefficient of $\sqrt2$.
Exploratory calculations (Observation~\ref{obs:window} below)
suggested that this is the general shape.  For an $S$-unit orbit choose
a representative $\xi=x+y\sqrt2\in\Z[\sqrt2]$ with $y\neq0$; its
reduction is $\xi(J)=x+2yJ$, and by \eqref{eq:lens-symmetries} the choice
of representative changes $\eta_J$ at most by sign.

\begin{conjecture}[Period law]
\label{conj:PL}
For every $S$-unit orbit with representative $\xi=x+y\sqrt2$, $y\neq0$,
the function $J\mapsto\eta_J\bigl(L(2J^2-1,\,x+2yJ)\bigr)$ is, for
$J\geq J_0(\xi)$, a quasi-polynomial in $J$ whose period divides $|y|$.
\end{conjecture}

\emph{Status.}  Conjecture~\ref{conj:PL} is proved for every orbit of
norm $\pm1$ (the silver powers, $m=0$) and of norm $\pm2$ ($|m|=1$) to
which it applies, that is with $y\neq0$; this excludes only the constant
orbit $q=\pm1$, where $\eta_J$ is cubic in $J$ (Lemma~\ref{lem:N0}):
Theorem~\ref{thm:PL} in Section~\ref{sec:period-law-proved}, with
Lemma~\ref{lem:S} and Theorem~\ref{thm:P} as its first ingredients.  For
$|m|\geq2$, for example the fourth family $2\delta=(1,2)$ and the orbit
$(1,-3)$ of $q\equiv4-4J$, it remains a conjecture, supported by the
held-out test below.

\emph{Held-out numerical test.}  Let $\mathcal W$ be the set of $64$ orbits
$(n,m)$ of $\delta^n\sqrt2^{\,m}$ with $0\leq n\leq5$, $|m|\leq9$
($m\geq0$ when $n=0$), $y\neq0$ and $|y|\leq41$; here
$|y|=2^{\lfloor|m|/2\rfloor}P_n$ for even $m$ and
$2^{\lfloor|m|/2\rfloor}a_n$ for odd $m$.  On $J\in[150,400]$ every orbit
of $\mathcal W$ was fitted by one polynomial per residue class of $J$
modulo $|y|$ (through the first four points of the class, verified on
the rest).  These fits were then used to predict $\eta_J$ for every orbit
of $\mathcal W$ and every $J\in[1401,1600]$, $12{,}800$ values, before
those rows were computed; all $12{,}800$ predictions were exact (only
the summary counts of this test were retained, not the individual
values).  This is finite evidence for Conjecture~\ref{conj:PL}; it does
not locate the onset $J_0(\xi)$, and nothing below uses it as a proof.

\begin{corollary}[Window]
\label{cor:U}
\begin{enumerate}[label=(\alph*)]
\item \textup{(Every rank; proved.)}  Among the $64$ orbits of
$\mathcal W$, exactly four vanish at every rank $J\geq2$: the orbits
$(0,1)$, $(1,-1)$, $(1,0)$, $(1,2)$ of Lemma~\ref{lem:sunit-four}, i.e.\
the four families of \cite[Thm.~7.1]{MahmoodP1Collection2026}.
\item \textup{(All large ranks.)}  For the orbits of $\mathcal W$ with
$|m|\leq1$ (norm $\pm1$ or $\pm2$), none but $(0,1)$, $(1,-1)$, $(1,0)$
vanishes at all sufficiently large ranks; this is proved
(Corollary~\ref{cor:R} below).  For the orbits with $|m|\geq2$, which
include the fourth family $(1,2)$, the corresponding statement is
observed on $J\in[150,400]\cup[1401,1600]$ and follows from
Conjecture~\ref{conj:PL} only together with the assumption that the
onset $J_0(\xi)$ of each such orbit lies below $150$.
\end{enumerate}
\end{corollary}

\begin{proof}
(a) The four orbits vanish at every $J\geq2$ by
\cite[Thm.~7.1]{MahmoodP1Collection2026}.  For each of the other $60$
orbits a single exact nonzero value refutes vanishing at every rank.
The following ledger gives one such witness $(J,q,\eta_J)$ per orbit
$(n,m)$, with $p=2J^2-1$ and $q$ the reduction of $\delta^n\sqrt2^{\,m}$
modulo $p$; every witness has $J\leq6$, and each value was cross-checked
against the independent evaluator of Appendix~\ref{app:verification}.
\begin{center}
\small
\begin{tabular}{@{}lrrr@{\qquad}lrrr@{\qquad}lrrr@{}}
\toprule
$(n,m)$ & $J$ & $q$ & $\eta_J$ & $(n,m)$ & $J$ & $q$ & $\eta_J$ & $(n,m)$ & $J$ & $q$ & $\eta_J$\\
\midrule
$(0,3)$ & 2 & 1 & $-2$ & $(2,-8)$ & 3 & 2 & $-2$ & $(3,-3)$ & 2 & 6 & $2$\\
$(0,5)$ & 4 & 1 & $-28$ & $(2,-7)$ & 2 & 1 & $-2$ & $(3,-2)$ & 5 & 4 & $-2$\\
$(0,7)$ & 4 & 2 & $-8$ & $(2,-6)$ & 4 & 14 & $2$ & $(3,-1)$ & 3 & 9 & $-2$\\
$(0,9)$ & 2 & 1 & $-2$ & $(2,-5)$ & 6 & 5 & $-2$ & $(3,0)$ & 2 & 6 & $2$\\
$(1,-9)$ & 4 & 10 & $2$ & $(2,-4)$ & 2 & 1 & $-2$ & $(3,1)$ & 3 & 1 & $-10$\\
$(1,-8)$ & 2 & 6 & $2$ & $(2,-3)$ & 5 & 33 & $-8$ & $(3,2)$ & 4 & 1 & $-28$\\
$(1,-7)$ & 3 & 9 & $-2$ & $(2,-2)$ & 3 & 16 & $10$ & $(3,3)$ & 2 & 6 & $2$\\
$(1,-6)$ & 5 & 32 & $4$ & $(2,-1)$ & 2 & 1 & $-2$ & $(3,4)$ & 4 & 2 & $-8$\\
$(1,-5)$ & 2 & 6 & $2$ & $(2,0)$ & 3 & 15 & $2$ & $(3,5)$ & 4 & 16 & $-8$\\
$(1,-4)$ & 4 & 10 & $2$ & $(2,1)$ & 4 & 28 & $2$ & $(3,6)$ & 2 & 6 & $2$\\
$(1,-3)$ & 3 & 2 & $-2$ & $(2,2)$ & 2 & 1 & $-2$ & $(4,-3)$ & 5 & 24 & $20$\\
$(1,-2)$ & 2 & 6 & $2$ & $(2,3)$ & 6 & 9 & $4$ & $(4,-2)$ & 2 & 1 & $-2$\\
$(1,1)$ & 2 & 6 & $2$ & $(2,4)$ & 3 & 9 & $-2$ & $(4,-1)$ & 5 & 48 & $60$\\
$(1,3)$ & 3 & 16 & $10$ & $(2,5)$ & 2 & 1 & $-2$ & $(4,0)$ & 4 & 20 & $2$\\
$(1,4)$ & 2 & 6 & $2$ & $(2,6)$ & 3 & 1 & $-10$ & $(4,1)$ & 2 & 1 & $-2$\\
$(1,5)$ & 3 & 15 & $2$ & $(2,7)$ & 6 & 36 & $-40$ & $(4,2)$ & 3 & 8 & $2$\\
$(1,6)$ & 4 & 10 & $2$ & $(2,8)$ & 2 & 1 & $-2$ & $(4,3)$ & 4 & 10 & $2$\\
$(1,7)$ & 2 & 6 & $2$ & $(3,-6)$ & 2 & 6 & $2$ & $(5,-1)$ & 2 & 6 & $2$\\
$(1,8)$ & 4 & 20 & $2$ & $(3,-5)$ & 3 & 15 & $2$ & $(5,0)$ & 5 & 37 & $-2$\\
$(1,9)$ & 3 & 9 & $-2$ & $(3,-4)$ & 5 & 2 & $-20$ & $(5,1)$ & 3 & 15 & $2$\\
\bottomrule
\end{tabular}
\end{center}
(b) The first statement is Corollary~\ref{cor:R}.  For the second, an
orbit that vanishes at all sufficiently large ranks has, under
Conjecture~\ref{conj:PL}, identically zero class polynomials from its
onset on; if the onset lies below $150$ these are the fitted
polynomials, which are not all zero for any of the other $60$ orbits.
\end{proof}

The every-rank statement (a) rests on exact witnesses and the four
all-rank proofs of \cite{MahmoodP1Collection2026}; it uses neither
Conjecture~\ref{conj:PL} nor the held-out test.

\begin{observation}[The zero set at $p=2J^2-1$; exploratory data]
\label{obs:window}
\begin{enumerate}[label=(\alph*)]
\item On $J\in[150,400]$ each orbit of $\mathcal W$ is a quasi-polynomial
in $J$ of degree at most $2$ on each residue class modulo $|y|$.  The
period-$2$ zero families in $\mathcal W$ are $(2,0)$
(Theorem~\ref{thm:C4}: even $J$), $(1,-3)$
(Remark~\ref{rem:C4-history}(iii): even $J$), $(1,-2)$ (odd $J$) and
$(0,3)$ (odd $J$).  Larger periods give scattered zero classes:
$\delta^5$ vanishes for $J\equiv0,\pm10\pmod{29}$ and $\delta^7$ for
$J\equiv0,\pm24,\pm46,\pm75\pmod{169}$.
\item The zero set is not confined to $S$-units.  At the ranks
$J=20,30,40$ ($p=799,1799,3199$) the numbers of units $q$ with
$\eta_J(L(p,q))=0$ are $120$, $180$ and $268$, among the $734$, $1{,}534$
and $2{,}734$ units $q\not\equiv\pm1$: about $6J$ zeros per rank.  About $41\%$ of
the units have no nonzero lattice point with $|u|,|v|\leq J-1$; between
$95\%$ and $100\%$ of the zeros at these ranks are of that kind, and for
them $\eta_J$ equals $-4Js+R_J$ alone.
\item What these counts do and do not show.  The zero set contains
infinitely many congruence families (Theorem~\ref{thm:P}); no $S$-unit
orbit of norm $\pm1$ or $\pm2$ other than the first three families of
\cite[Thm.~7.1]{MahmoodP1Collection2026} vanishes at every rank
(Corollary~\ref{cor:R}); and in $\mathcal W$ exactly the four families
vanish at every rank (Corollary~\ref{cor:U}(a)).  Whether the full zero
set at $p=2J^2-1$ admits a finite parametrized description is open: an
infinite indexed ladder can itself have one finite defining formula,
and the counts at three ranks do not decide this.
\end{enumerate}
\end{observation}

\begin{remark}[A refuted ``zero-term principle'']
\label{rem:zero-term}
One might guess that on a zero class each of $-4Js+R_J$, $\Delta_J$ and
$L_J$ vanishes separately, as in slice (i) of Theorem~\ref{thm:P}.  This
is false: in $52$ of $4{,}383$ zero-class rows examined, all in the
orbits $(4,\pm2)$, the term $-4Js+R_J$ cancels against $L_J$ instead.
Zeros of $\eta_J$ can arise from cancellation between the
Dedekind--residue part and the lattice part.
\end{remark}

\begin{remark}[Narrowing of Observation~\ref{obs:rho-statistics}]
\label{rem:narrowing}
Observation~\ref{obs:rho-statistics}(b) describes statistics in the
coordinate $\rho$ at primes taken one silver rank each.  It does not test
congruence classes of ranks along $p=2J^2-1$: in the coordinate $q$ the
silver-power orbits carry structured zeros on such classes
(Theorem~\ref{thm:P}).  Thus those finite coordinate statistics do not
justify describing these zeros as accidental.  No structural claim is
made about the remaining sporadic zeros, and
Observation~\ref{obs:rho-statistics} remains a description of its stated
ranges.
\end{remark}

\section{The period law in the unit and norm-two sectors}
\label{sec:period-law-proved}

Conjecture~\ref{conj:PL} is a theorem for every nonconstant $S$-unit
orbit of norm $\pm1$ or $\pm2$, that is every such orbit with $y\neq0$.
Concretely, with $a=a_n$, $b=P_n$, $\nu=\nu_n$ as in
Section~\ref{sec:pell}, the orbits of Theorem~\ref{thm:PL} are
\begin{itemize}
\item the \emph{unit sector} $\mathsf S_0$: $q=a_n+2P_nJ=q_n$, the
reduction of $\delta^n$, $n\geq1$, of period $P_n$;
\item the \emph{norm-two sector} $\mathsf S_1$: $q=a_nJ+\sigma P_n$
with $\sigma=\pm1$, $n\geq2$, of period $a_n$; here $\sigma=+1$ is the
reduction of $\delta^n/\sqrt2$, the orbit $(n,-1)$, and $\sigma=-1$ is,
up to the sign $\nu$, the reduction of $\delta^{-n}/\sqrt2$, the orbit
$(n,1)$ of $\delta^n\sqrt2$;
\end{itemize}
and the small nonconstant orbits $q=J$ and $q=J\pm1$ ($n=0,1$ in
$\mathsf S_1$, period $1$) are covered by Theorem~\ref{thm:C3},
\cite[Thm.~7.1]{MahmoodP1Collection2026} and Lemma~\ref{lem:N}, where
$\eta_J$ is $0$, $0$ and $-2$.  The constant orbit $q=\pm1$ ($n=0$ in
$\mathsf S_0$, $y=0$) lies outside Conjecture~\ref{conj:PL}; there
$\eta_J(L(2J^2-1,1))=-J(J-1)(2J-1)/3$ is cubic in $J$
(Lemma~\ref{lem:N0}), and no statement below about linear class
polynomials applies to it.  An orbit $(n,m)$ has norm
$\pm1$ iff $m=0$ and contains an element of norm $\pm2$ iff $|m|=1$; the
fourth family $2\delta=(1,2)$ has norm $-4$ and lies outside this
section.  The proof has four ingredients: Lemma~\ref{lem:S} and its
norm-two analogue make the Dedekind term a rational function of $k$ on
each residue class; a breakpoint lemma does the same for the residue
term; Theorem~\ref{thm:U} bounds the denominators of $\eta_J$, which
forces a rational function to be a polynomial; and a box form of the two
lattice terms of \eqref{eq:master} shows that they are eventually linear
in $k$.  Throughout, $\Lambda_q=\{(u,v)\in\Z^2:\ v+qu\equiv0\pmod p\}$,
$\|(u,v)\|=\max(|u|,|v|)$, and $\langle x\rangle$ is the least
non-negative residue modulo $p$.

\begin{lemma}[Box form of the lattice terms]
\label{lem:box}
For every $p\geq2$, every unit $q$ modulo $p$ and every $J\geq1$,
\begin{equation}
 -2\Delta_J-L_J=\sum_{\substack{(u,v)\in\Lambda_q\setminus\{0\}\\
 \|(u,v)\|\leq J-1}}\sgn(uv)\,\bigl(J-\|(u,v)\|\bigr)=:\mathcal B_J(q).
 \label{eq:box}
\end{equation}
Hence $\eta_J(L(p,q))=\bigl[-4J\,s(q,p)+R_J\bigr]+\mathcal B_J(q)$.  This
is a rewriting of the two lattice terms of \eqref{eq:master}, not a new
invariant.
\end{lemma}

\begin{proof}
\emph{The low-mode term.}  By the substitution of the proof of
Theorem~\ref{thm:C4}, $w_L=u+v$, $w_R=v-u$, the pairs counted by
$I_{J-m,J+m}$ are the lattice points $(u,v)\in\Lambda_q$ with
$|u+v|\leq J-m-1$, $|v-u|\leq J+m-1$ and $u+v\equiv J-m-1\pmod2$, and
those counted by $I_{J+m,J-m}$ are the same with the two bounds
exchanged; every such point has $\|(u,v)\|\leq J-1$.  Fix a nonzero
$(u,v)\in\Lambda_q$ with $M:=\|(u,v)\|\leq J-1$, put
$m_0:=\min(|u|,|v|)$, $r=u+v$, $t=v-u$ (so $r\equiv t\pmod2$), and
count the $m\in[1,J-1]$ for which the point is counted.  If $|r|\leq J$
and $|t|\leq J$, the two counts are $\lfloor(J-|r|)/2\rfloor$ and
$\lfloor(J-|t|)/2\rfloor$, whose difference is
$(|t|-|r|)/2=-\sgn(uv)\,m_0$.  Otherwise exactly one of the two counts
is nonzero and equals $J-M$, the count with the sign
$-\sgn(uv)$.  Hence the point contributes $-\sgn(uv)\min(m_0,\,J-M)$
to $L_J$; points with $uv=0$ contribute nothing.

\emph{The triangular sum.}  A pair $(a,b)$ with $aq\equiv b$ is the
point $(a,-b)\in\Lambda_q$ and a pair with $aq\equiv-b$ is the point
$(a,b)$, with weight $J-|u|-|v|$ when $|u|+|v|\leq J-1$; so, using the
symmetry $(u,v)\mapsto(-u,-v)$ of $\Lambda_q$ and of $\sgn(uv)$,
$-2\Delta_J=\sum\sgn(uv)\,(J-|u|-|v|)_+$ over the nonzero points of
$\Lambda_q$.

\emph{Sum.}  If $m_0+M\leq J-1$ the two contributions add to
$(J-M-m_0)+m_0=J-M$; otherwise $(J-|u|-|v|)_+=0$ and
$\min(m_0,J-M)=J-M$.
\end{proof}

\begin{lemma}[Norm-two Dedekind reduction]
\label{lem:Sprime}
Let $n\geq1$, $J\geq2$, $p=2J^2-1$, $\sigma=\pm1$, $q=aJ+\sigma b$ and
$t=aJ-\sigma b$.  Then $a^2p=2qt-\nu$, hence $q^*\equiv2\nu t\pmod p$, and
\begin{equation}
 s(q,p)=\mathcal T(q,p)+\nu\,\mathcal T(a^2,q)-\nu\,s\bigl(q\bmod a^2,\,a^2\bigr),
 \label{eq:lemma-Sprime}
\end{equation}
where $q\bmod a^2=(aJ+\sigma b)\bmod a^2$ depends only on $J\bmod a$.
\end{lemma}

\begin{proof}
$2qt=2(a^2J^2-b^2)=a^2(p+1)-2b^2=a^2p+\nu$ by (F1).  So $2qt\equiv\nu$
modulo $p$, which gives $q^*$, and $a^2p\equiv-\nu\pmod q$, so
$p\equiv-\nu a^{-2}\pmod q$; here $\gcd(a,q)=\gcd(a,b)=1$.  Hence
$s(p,q)=-\nu\,s(a^2,q)$, and two applications of
\eqref{eq:reciprocity} give \eqref{eq:lemma-Sprime}.
\end{proof}

\begin{lemma}[Breakpoints]
\label{lem:K}
Let $k\geq1$.
\begin{enumerate}[label=(\alph*)]
\item \textup{($\mathsf S_0$.)}  Let $J=bk+j$ with $0\leq j<b$,
$q=q_n$, $t=2bJ-a$, $d_j=2bj-a$ and $q_j=a+2bj$.  For $1\leq c\leq b$,
$\lceil cp/q\rceil=ck+e_c$ with $e_c=\lceil c\,d_j/(2b^2)\rceil$, or
$c\,d_j/(2b^2)+\ind_{\nu=1}$ when that quotient is an integer; and
$\lceil cp/t\rceil=ck+e'_c$, the same with $q_j$ in place of $d_j$.
\item \textup{($\mathsf S_1$.)}  Let $J=ak+j$ with $0\leq j<a$,
$q=aJ+\sigma b$, $t=aJ-\sigma b$.  Then
$\lceil cp/q\rceil=2ck+\lceil2c(aj-\sigma b)/a^2\rceil$ for
$1\leq c\leq(a+1)/2$ and
$\lceil cp/(2t)\rceil=ck+\lceil c(aj+\sigma b)/a^2\rceil$ for
$1\leq c\leq a$, both ceilings being of non-integers.
\item In both sectors $q^*\equiv-\nu t$ ($\mathsf S_0$) or $2\nu t$
($\mathsf S_1$).  The number of wraps $\lfloor\alpha x/p\rfloor$ of a
residue $\langle\alpha x\rangle$, $1\leq\alpha\leq J-1$, equals the number
of the breakpoints $\lceil cp/x\rceil$ that are $\leq\alpha$.  The
breakpoints entering the layer-cake sums below ($c\leq b-1$ in
$\mathsf S_0$; $c\leq(a-1)/2$ for $q$ and $c\leq a-1$ for $2t$ in
$\mathsf S_1$) are all $\leq J+1$, and in $\mathsf S_1$ all $\leq J$,
whenever $J\geq P$ (the period).  The $\mathsf S_1$ breakpoint
$c=(a+1)/2$ of $q$ is listed only to show that it is $\geq J$; it is
inactive and can exceed $J+1$ (for $a=3$, $b=2$, $\sigma=-1$, $J=3$ it is
$\lceil34/7\rceil=5$).
\end{enumerate}
\end{lemma}

\begin{proof}
(a) From $M_np=qt+\nu$ \eqref{eq:F5} and $t=M_nk+d_j$, $q=M_nk+q_j$:
$cp/q=ck+c\,d_j/M_n+c\nu/(M_nq)$ and $cp/t=ck+c\,q_j/M_n+c\nu/(M_nt)$.
The last terms have absolute value $<1/M_n$ because $c<q$ and $c<t$, so
the ceiling is decided by $c\,d_j/M_n$ (resp.\ $c\,q_j/M_n$), or, when
that is an integer, by the sign of $\nu$.  Also $\lceil bp/q\rceil=J$
and $\lceil bp/t\rceil=J+1$, because $0<a/(2b)<1$.
(b) From $a^2p=2qt-\nu$ (Lemma~\ref{lem:Sprime}) and $t=a^2k+aj-\sigma b$,
$q=a^2k+aj+\sigma b$: $cp/q=2ck+2c(aj-\sigma b)/a^2-c\nu/(a^2q)$ and
$cp/(2t)=ck+c(aj+\sigma b)/a^2-c\nu/(2a^2t)$.  In the stated ranges
$a^2$ divides neither $2c(aj-\sigma b)$ nor $c(aj+\sigma b)$, because
every prime $\ell\mid a$ has $\ell\nmid b$, hence
$\ell\nmid(aj\mp\sigma b)$; so the ceilings are of non-integers, and
the corrections are below $1/a^2$ in absolute value.
(c) The inverses are \eqref{eq:F5} and Lemma~\ref{lem:Sprime}; the wrap
count is the definition of the breakpoints.  In $\mathsf S_1$ the
breakpoint $c=(a+1)/2$ of $q$ is at least $(a+1)k+j-1>J-1$, and
$\lceil ap/(2t)\rceil\geq J$; in $\mathsf S_0$ the last breakpoints are
$J$ and $J+1$ by (a), and they bound every earlier active breakpoint by
monotonicity in $c$.  For the active $\mathsf S_1$ breakpoints we use the
following identities: with $a$ odd,
\[
 Jq-\tfrac{a-1}2\,p=J^2+\sigma bJ+\tfrac{a-1}2,\qquad
 J(2t)-(a-1)\,p=2J^2-2\sigma bJ+a-1,
\]
both positive for $J\geq a>b$.  Hence
$\lceil\frac{a-1}2\,p/q\rceil\leq J$ and $\lceil(a-1)p/(2t)\rceil\leq J$,
and by monotonicity in $c$ every active breakpoint is $\leq J$.  The
layer-cake form is used below only for $k\geq1$, that is for $J\geq P$,
which in $\mathsf S_1$ is $J\geq a$.
\end{proof}

The finite breakpoint checks, on every class $j<40$ with $k\leq5$ of
the $24$ families of Observation~\ref{obs:zero-classes}, are reported
separately in Appendix~\ref{app:verification}.

\begin{theorem}[The period law in $\mathsf S_0$ and $\mathsf S_1$]
\label{thm:PL}
Let the family be $\mathsf S_0$ ($n\geq1$) or $\mathsf S_1$ ($n\geq2$,
$\sigma=\pm1$), with period $P=P_n$ ($\mathsf S_0$) or $P=a_n$
($\mathsf S_1$).  For each residue $j$ modulo $P$ there is a polynomial
$\mathcal Q_j$ of degree at most $1$ and a threshold $k^*_j\geq1$, computed
from linear data of the class, such that
\begin{equation}
 \eta_J(L(p,q))=\mathcal Q_j(k)\qquad\text{for every }J=Pk+j\text{ with }k\geq k^*_j .
 \label{eq:PL}
\end{equation}
A class is a zero class ($\mathcal Q_j\equiv0$) iff $\mathcal Q_j$
vanishes at $k^*_j$ and $k^*_j+1$.  The complete zero set of
the orbit is the union of the zero classes (for $k\geq k^*_j$), the
isolated zeros (integer roots $k\geq k^*_j$ of the nonzero
$\mathcal Q_j$, at most one per class), and the finitely many zeros with
$k<k^*_j$, which are checked directly.  A nonzero class polynomial can
vanish at an isolated rank, so the zero set need not be a union of whole
residue classes beyond the thresholds.
\end{theorem}

\begin{proof}
Write $\eta_J=f(k)+\mathcal B_J$ with $f=-4Js+R_J$ (Lemma~\ref{lem:box}),
and fix the class $j$.

\emph{(1) The bulk $f$ is a rational function of $k$ for every $k\geq1$.}
The Dedekind term is given by Lemma~\ref{lem:S} ($\mathsf S_0$) or
Lemma~\ref{lem:Sprime} ($\mathsf S_1$); its only arithmetic input is a
Dedekind sum of fixed modulus ($2b^2$ or $a^2$) at an argument fixed by
$j$, and the rest is the rational function $\mathcal T$.  The residue
term is $\frac2p\sum_\alpha(J-\alpha)\rho_p(\alpha q)$ plus the inverse
part, which is $-\nu\frac2p\sum_\alpha(J-\alpha)\rho_p(\alpha t)$ in
$\mathsf S_0$ and $\nu\frac2p\sum_\alpha(J-\alpha)\rho_p(2\alpha t)$ in
$\mathsf S_1$ (Lemma~\ref{lem:K}(c)).  Here
$\rho_p(\alpha x)=p\,(1+2c_x(\alpha))-2\alpha x$ with $c_x(\alpha)$ the
number of wraps, and by Lemma~\ref{lem:K} and the layer-cake identity
$\sum_\alpha(J-\alpha)c_x(\alpha)=\sum_{\beta}(J-\beta)(J-\beta+1)/2$
over the active breakpoints $\beta$ of Lemma~\ref{lem:K}(c) (all
$\leq J+1$, so a breakpoint equal to $J$ or $J+1$ contributes $0$),
$p\,R_J$ is a polynomial in $k$.

\emph{(2) $f(k)\in\frac23\Z$ for every $k\geq1$.}  By
Theorem~\ref{thm:U} at $p=2J^2-1$, $\eta_J\in-4Jp\,s(q,p)+2\Z$, and
$6p\,s(q,p)\in\Z$ by Lemma~\ref{lem:dedekind-facts}, so
$\eta_J\in\frac23\Z$; and $\mathcal B_J\in2\Z$, because its summand is
invariant under $(u,v)\mapsto(-u,-v)$ and $(0,0)$ is excluded, so the
points pair up.  (An integer $\mathcal B_J$ would already give
$f\in\frac13\Z$, which also suffices below.)

\emph{(3) Rational to polynomial.}  A rational function
$g\in\Q(k)$ with $g(k)\in\frac1D\Z$ for all large integers $k$ is a
polynomial: write $g=S+R/Q$ with $\deg R<\deg Q$; then
$D'\,R(k)/Q(k)\in\Z$ for large $k$ and tends to $0$, so $R\equiv0$.
Growth: $s=O(J)$ by Lemma~\ref{lem:S} or \ref{lem:Sprime}, so
$f=O(k^2)$ and $\deg f\leq2$.

\emph{(3$'$) The $k^2$ coefficient of $f$ vanishes.}  Write
$\rho_p(\alpha x)=p(1+2c_x(\alpha))-2\alpha x$ with $c_x(\alpha)$ the
number of wraps.  Then $R_J$ is the sum of a constant part
$2(1\mp\nu)\sum_\alpha(J-\alpha)$ (upper sign $\mathsf S_0$), the
layer-cake part $4\sum_\alpha(J-\alpha)\bigl(c_q(\alpha)\mp\nu\,c_x(\alpha)\bigr)$
with $x=t$ ($\mathsf S_0$) or $x=2t$ and the lower sign
($\mathsf S_1$), and the part $-\frac4p\sum_\alpha(J-\alpha)\alpha\,(q-\nu t)$,
resp.\ $(q+2\nu t)$, with $\sum_\alpha(J-\alpha)\alpha=(J^3-J)/6$.  The
breakpoints are linear in $k$ with bounded offsets (Lemma~\ref{lem:K}),
and the Dedekind lookups are bounded, so only leading slopes enter the
$k^2$ coefficient $\lim f/k^2$.  Part by part it is
\begin{center}
\small
\begin{tabular}{@{}lccccc@{}}
\toprule
& $-4J\cdot$(rational part of $s$) & constant & layer cake & $\alpha$-part & sum\\
\midrule
$\mathsf S_0$, $\nu=+1$ & $0$ & $0$ & $0$ & $0$ & $0$\\
$\mathsf S_0$, $\nu=-1$ & $-\frac{2b}3$ & $2b^2$ & $\frac23(b-1)b(2b-1)$ & $-\frac43b^3$ & $0$\\
$\mathsf S_1$, $\nu=+1$ & $-a$ & $2a^2$ & $a(a-1)^2$ & $-a^3$ & $0$\\
$\mathsf S_1$, $\nu=-1$ & $-\frac a3$ & $0$ & $-\frac13(a-1)a(a+1)$ & $\frac{a^3}3$ & $0$\\
\bottomrule
\end{tabular}
\end{center}
independently of $\sigma$, so $\deg f\leq1$.  (The values were derived by
hand and certified symbolically by sympy and by the local Wolfram
Engine~14.3, with identical part values; Remark~\ref{rem:period-certificates}.)

\emph{(4) The box $\mathcal B_J$ is linear in $k$ for $k\geq k^*_j$.}
Within one wrap interval $\langle\alpha q\rangle$ increases by $q>J-1$,
so a lattice point of the box with $u=\alpha\geq1$ sits at an interval
start ($\langle\alpha q\rangle\leq J-1$, $v=-\langle\alpha q\rangle$) or
at an interval end ($p-\langle\alpha q\rangle\leq J-1$,
$v=p-\langle\alpha q\rangle$); negative $u$ follow by symmetry, and
$u=0$ gives $v=0$.  At the starts and at the regular ends
($\alpha=\beta-1$ for a breakpoint $\beta\leq J$) both $u$ and $|v|$ are
linear in $k$: $\langle\alpha_cq\rangle=\alpha_cq-cp$ has no $k^2$ term,
and $p-\langle(\beta-1)q\rangle=q-\langle\beta q\rangle$.  The last
element $u=J-1$ is a regular end, except in $\mathsf S_1$ when it lies
in the interval $c=(a-1)/2$ truncated at $J-1$; there
$p-\langle(J-1)q\rangle=\tfrac{a+1}2p-(J-1)q=J^2+(a-\sigma b)J+\sigma b-\tfrac{a+1}2$,
which exceeds $J-1$ for $J\geq a$ because $a^2>b+(a-1)/2$, so it is
never in the box (if that breakpoint exceeds $J-1$, which happens only
at $k=1$, it equals $J$ and $J-1$ is a regular end).  Now $k^*_j\geq1$ is
the first $k\geq1$ beyond every crossing of the linear conditions $|v|\leq J-1$,
$u\leq J-1$ and $u\lessgtr|v|$; each linear form is fixed by three exact
values, and collinearity certifies that its $k^2$ term is absent.  For
$k\geq k^*_j$, $\mathcal B_J$ is a fixed signed sum of terms
$J-\max(u,|v|)$, hence linear in $k$.

\emph{(5)} $\eta_J=f+\mathcal B_J$ is a polynomial of degree at most $1$
for $k\geq k^*_j$; two values determine it.
\end{proof}

\emph{The decision procedure.}  For a family and a class $j$, compute
the threshold $k^*_j$ from the linear data of step~(4), evaluate
$\eta_J$ exactly at $k=k^*_j,k^*_j+1,k^*_j+2$, and read off
$\mathcal Q_j$; the class is a zero class iff the values vanish (two
suffice by the degree bound; the recorded computations used three).
For the complete zero set, also solve each nonzero $\mathcal Q_j$ for
integer roots $k\geq k^*_j$ and evaluate the finitely many ranks with
$k<k^*_j$ directly.  The zero classes for $\mathsf S_0$, $n\leq10$, and
$\mathsf S_1$, $n\leq8$, are in Observation~\ref{obs:zero-classes}.
Across these $24$ families there is exactly one isolated zero beyond the
thresholds: $\mathsf S_1$, $n=7$, $\sigma=+1$, class $j=190$, where
$\mathcal Q_{190}(k)=2-2k$ vanishes at $k=1$, $J=429$ (and
$\eta=-2,-4$ at $J=668,907$).  There are also $363$ zeros below the
thresholds, mostly at the smallest rank of a class.  The isolated zero
and all $364$ zeros were re-evaluated exactly;
the complete sets are retained with the sources
(\path{results/LC14B_ZERO_SETS.json} of \cite{LCrecord}).

\begin{corollary}[Eventual periodicity of the zero set]
\label{cor:eventual}
In each family of Theorem~\ref{thm:PL}, with period $P$, put
\[
 J_1=\max_j\Bigl(P\,(k^*_j+1)+j,\ \max_r\bigl(P\,(r+1)+j\bigr)\Bigr),
\]
where $r$ runs over the integer roots of the nonzero $\mathcal Q_j$ (the
inner maximum is omitted when there are no such roots).  For
$J\geq J_1$, $\eta_J(L(2J^2-1,q))=0$ iff $J\bmod P$ is a zero class.
Hence the zero set of the orbit is eventually periodic, and its natural
density is (number of zero classes)$/P$.
\end{corollary}

\begin{proof}
For $J=Pk+j\geq J_1$ one has $k>k^*_j$ and $k$ exceeds every root of
$\mathcal Q_j$, so $\eta_J=\mathcal Q_j(k)$ vanishes iff
$\mathcal Q_j\equiv0$.
\end{proof}

The periodic onset need not equal a class threshold:
the isolated zero $J=429$ lies past its class threshold but below $J_1$.
Regular class laws and one-off cancellations coexist.  For example, the
densities of the zero sets of $\delta^3,\ldots,\delta^9$ are $1/5$,
$2/12$, $3/29$, $10/70$, $7/169$, $18/408$ and $3/985$, read off from
Observation~\ref{obs:zero-classes}.

\begin{theorem}[The norm-two slice $J\equiv0\pmod{a_n}$]
\label{thm:Q}
For every $n\geq2$, $k\geq1$ and $J=a_nk$, $p=2J^2-1$,
\[
 \eta_J\bigl(L(p,\,a_nJ-P_n)\bigr)=
 \begin{cases}-P_{n-1}^2\,k,&n\text{ odd},\\[2pt]
 \tfrac13P_nP_{n-2}\,k,&n\text{ even},\end{cases}
 \qquad
 \eta_J\bigl(L(p,\,a_nJ+P_n)\bigr)=-\eta_J\bigl(L(p,\,a_nJ-P_n)\bigr).
\]
It is zero only for $n=2$.
\end{theorem}

\begin{proof}
This is the class $j=0$ of $\mathsf S_1$.  By Lemma~\ref{lem:Sprime},
since $q\bmod a^2\equiv\sigma b$ and $a^2\equiv\nu\pmod b$ by (F1),
$s(q,p)=\mathcal T(q,p)+\nu\mathcal T(a^2,q)-\nu\sigma\mathcal T(b,a^2)+\sigma s(1,b)$.
Lemma~\ref{lem:K}(b) gives the breakpoints $2ck+\ind_{\sigma=-1}$ for
$q$, $1\leq c\leq(a-1)/2$, and $ck+\ind_{\sigma=1}$ for $2t$,
$1\leq c\leq a-1$.  The box is empty for every $k\geq1$: at the interval
starts $\langle\alpha_cq\rangle=k(2c\sigma b+\ind_{\sigma=-1}a^2)+(\ind_{\sigma=-1}\sigma b+c)$
exceeds $J-1$; at the regular ends $p-\langle\alpha q\rangle$ equals
$(c+1)(2bk-1)$ for $\sigma=-1$ and $k(a^2-2(c+1)b)+(b-c-1)$ for
$\sigma=+1$, both exceeding $J-1$ by (F4) ($b\leq a<2b$); and the final
end is excluded as in step~(4) of Theorem~\ref{thm:PL}.  So
$\eta_J=f=-4Js+R_J$, and the Faulhaber sums give the stated values
(certificate, Remark~\ref{rem:period-certificates}).  For $n\geq3$ the
values are nonzero since $P_{n-1},P_{n-2}\geq1$; for $n=2$, $P_0=0$.
\end{proof}

\begin{lemma}
\label{lem:N}
$\eta_J(L(2J^2-1,\,J-1))=-2$ for every $J\geq2$.
\end{lemma}

\begin{proof}
$p=(J-1)(2J+2)+1$, so $s(J-1,p)=\mathcal T(J-1,p)-s(1,J-1)$ and
$q^*\equiv-(2J+2)$.  No residue wraps, as $\alpha(2J+2)\leq p-1$ for
$\alpha\leq J-1$.  The box is $\{\pm(1,-(J-1))\}$, giving
$\mathcal B_J=-2$, and the bulk $-4Js+R_J$ vanishes identically
(certificate, Remark~\ref{rem:period-certificates}).
\end{proof}

\begin{lemma}
\label{lem:N0}
$\eta_J(L(2J^2-1,\,1))=-J(J-1)(2J-1)/3$ for every $J\geq2$.
\end{lemma}

\begin{proof}
$s(1,p)=(p-1)(p-2)/(12p)$ and $\rho_p(\alpha)=p-2\alpha$.  The box is
$\{(u,-u):\ 1\leq|u|\leq J-1\}$, since $p>2J-2$, so
$\mathcal B_J=-J(J-1)$.  Summing gives the value (certificate,
Remark~\ref{rem:period-certificates}).
\end{proof}

\begin{corollary}[Rank-uniform zeros of norm $\pm1$ and $\pm2$]
\label{cor:R}
The $S$-unit orbits of norm $\pm1$ or $\pm2$ are $\pm1$,
$\pm\delta^{\pm n}$ and the norm-two orbits $a_nJ\pm P_n$ ($n\geq0$).
Those for which $\eta_J$ vanishes at every rank $J\geq2$, or at all
sufficiently large ranks, are exactly the first three families of
\cite[Thm.~7.1]{MahmoodP1Collection2026}: $q\equiv J$, $J+1$, $2J-1$, up
to orbit.  The fourth family $2\delta$ has norm $-4$ and is not covered.
\end{corollary}

\begin{proof}
\emph{Norm $\pm1$.}  $q=\pm1$ never vanishes (Lemma~\ref{lem:N0});
$\delta$ is the third family; $\delta^n$ with $n\geq2$ is nonzero on
$J\equiv1\pmod{P_n}$ (Theorem~\ref{thm:P}(ii)).
\emph{Norm $\pm2$, $n\geq3$.}  These orbits are nonzero on
$J\equiv0\pmod{a_n}$ (Theorem~\ref{thm:Q}).
\emph{$n=2$, i.e.\ $q=3J\mp2$.}  Both orbits are nonzero on the class
$J\equiv1\pmod3$, where the procedure's certified polynomial is
$\mathcal Q(k)=2k$ for both signs.
\emph{$n=1$.}  $q=J+1$ is the second family; $q=J-1$ gives $-2$
(Lemma~\ref{lem:N}).
\emph{$n=0$.}  $q=J$ is the first family.
\end{proof}

\begin{observation}[Zero classes in the proved sectors]
\label{obs:zero-classes}
The decision procedure gives the complete zero-class sets below
($J\equiv j$ modulo the period); the entries for $\delta^{10}$ and for
$n=8$ in $\mathsf S_1$ were never computed before the procedure was
fixed.  For $\delta^{10}$ the $58$ classes are enumerated in the
supplementary JSON cited after the table; only their number is printed
in the table.
\begin{center}
\small
\begin{tabular}{@{}lcL{0.58\textwidth}@{}}
\toprule
orbit & period & zero classes\\
\midrule
$\delta$ & $1$ & all $J\geq2$ (third family)\\
$\delta^2$ & $2$ & $\{0\}$ (Theorem~\ref{thm:C4})\\
$\delta^3$ & $5$ & $\{0\}$\\
$\delta^4$ & $12$ & $\{3,9\}$\\
$\delta^5$ & $29$ & $\{0,10,19\}$\\
$\delta^6$ & $70$ & $\{7,10,13,14,20,50,56,57,60,63\}$\\
$\delta^7$ & $169$ & $\{0,24,46,75,94,123,145\}$\\
$\delta^8$ & $408$ & $\{17,24,34,48,72,85,96,109,119,\allowbreak
289,299,312,323,336,360,374,384,391\}$\\
$\delta^9$ & $985$ & $\{0,193,792\}$\\
$\delta^{10}$ & $2378$ & $58$ classes\\
\midrule
$a_nJ-P_n$ / $a_nJ+P_n$, $n=2$ & $3$ & $\{0\}$ / $\{0\}$\\
$n=3$ & $7$ & $\{2\}$ / $\{5\}$\\
$n=4$ & $17$ & $\{4,13,14\}$ / $\{3,4,13\}$\\
$n=5$ & $41$ & $\{8,12,27\}$ / $\{14,29,33\}$\\
$n=6$ & $99$ & $\{20,21,28,31,71,78,79\}$ / $\{20,21,28,68,71,78,79\}$\\
$n=7$ & $239$ & $\{41,46,70,75,113,193\}$ / $\{46,126,164,169,193,198\}$\\
$n=8$ & $577$ & $\{32,119,120,122,168,409,\allowbreak410,439,457,458,547\}$ /
$\{30,119,120,138,167,168,\allowbreak409,455,457,458,545\}$\\
\bottomrule
\end{tabular}
\end{center}
The full list of $58$ classes for $\delta^{10}$ is included in
\path{results/LC14B_ZERO_SETS.json} of the supplement \cite{LCrecord}.
Across these $24$ families the largest threshold $k^*_j$ is $2$ and
every $\mathcal Q_j$ has degree at most $1$, as Theorem~\ref{thm:PL}
now proves.  \emph{Observed, not proved:} the zero sets of $\delta^n$
are symmetric under $j\mapsto-j$, and the two signs of the norm-two
orbits have mirror-image zero sets.
\end{observation}

\begin{remark}[Symbolic certificates]
\label{rem:period-certificates}
The Faulhaber-sum evaluations behind Theorem~\ref{thm:Q} (four cases
$(\nu,\sigma)$), the vanishing bulk of Lemma~\ref{lem:N} and the sum of
Lemma~\ref{lem:N0} were certified symbolically by sympy (exact
rationals, polynomial remainder modulo $a^2-2b^2-\nu$, no numerical
simplification), together with a re-run of the identities of
Remark~\ref{rem:faulhaber-certificate}; seven identity sets, all true.
The same identities were independently checked in Wolfram
Engine~14.3.  Theorem~\ref{thm:Q} in all four cases and
Lemmas~\ref{lem:N} and~\ref{lem:N0} agree in both routes.  The
certificates of Theorem~\ref{thm:P} remain those of
Remark~\ref{rem:faulhaber-certificate}, sympy and Wolfram Engine~14.3.
The leading-coefficient table of step~(3$'$) of Theorem~\ref{thm:PL}
was certified separately by sympy and by Wolfram
Engine~14.3: each part's $k^2$ coefficient was computed as
$\lim_{k\to\infty}(\text{part})/k^2$ from the exact expressions, with
the layer-cake parts as symbolic sums over $c$; the two routes give
identical part values and total $0$ in all six cases $(\mathsf S_0,\nu)$
and $(\mathsf S_1,\nu,\sigma)$.
Every identity is also confirmed by exact evaluation on every row of
Appendix~\ref{app:verification}.
\end{remark}

\begin{remark}[Verification]
\label{rem:period-verification}
Lemma~\ref{lem:box} was checked against the literal $\Delta_J$ and $L_J$
on $181$ random triples $(p,q,J)$ with arbitrary moduli $2\leq p\leq150$
and $J\leq40$, and the mutated weight $J-|u|-|v|$ fails.
Theorem~\ref{thm:Q} was checked on $70$ exploratory rows ($n\leq8$,
$2\leq k\leq6$, both signs) and $82$ rows fixed before computation
($n\leq12$ with $k=7,8,9$ and $13\leq n\leq16$ with $k=1,2$, both
signs), the largest being $n=16$, $J=1{,}331{,}714$ with
$\eta_J(L(p,a_{16}J-P_{16}))=25{,}356{,}500{,}416=\frac13P_{16}P_{14}\cdot2$;
also $n=15$, $J=551{,}614$: $\eta_J=-13{,}051{,}463{,}048=-P_{14}^2\cdot2$;
the sign-flipped statement fails on every odd-$n$ held-out row.
Lemmas~\ref{lem:N} and \ref{lem:N0} were checked at $J=401,\ldots,410$
and at $J=10{,}000$, resp.\ $J=5{,}000$.  For Theorem~\ref{thm:PL} the
procedure's polynomial $\mathcal Q_j$ of every class of the $24$
families predicted $\eta_J$ at $k^*_j+3$ and $k^*_j+4$ (and, for
$\delta^n$ with $n\leq8$ and the norm-two orbits with $n\leq6$, at
$3k^*_j+11$) before those rows were computed: $13{,}080$ predictions,
none missed.  A fit with the wrong period $P_n+1$ mispredicts $6$, $13$,
$30$, $71$ and $170$ classes for $n=3,\ldots,7$.
\end{remark}

\begin{corollary}[Unbounded ambiguity of a vanishing eta invariant]
\label{cor:ambiguity}
For every $K\geq1$ there are a rank $J$ and the level $p=2J^2-1$ with
$K$ chiral lens spaces $L(p,q_1),\ldots,L(p,q_K)$ whose orbits
$\{\pm q_i^{\pm1}\}$ are pairwise distinct (so no two are homeomorphic,
even after orientation reversal), all with $\eta_J(L(p,q_i))=0$ and
none rank-$J$ balanced.  One can take the $q_i$ to be the reductions of
$\delta^{n_i}$ for any $K$ distinct odd $n_i\geq3$, and $J$ any
sufficiently large multiple of
$\operatorname{lcm}(P_{n_1},\ldots,P_{n_K})$.
\end{corollary}

\begin{proof}
Fix distinct odd $n_1,\ldots,n_K\geq3$, put
$\Lambda=\operatorname{lcm}(P_{n_1},\ldots,P_{n_K})$ and let $J$ be a
multiple of $\Lambda$.  Then $P_{n_i}\mid J$ and Theorem~\ref{thm:P}(i)
gives $\eta_J(L(p,q_{n_i}))=0$ for every $i$.  Since
\eqref{eq:sqrt2-reduction} is a ring homomorphism,
$q_nq_m\equiv q_{n+m}\pmod p$.  Hence each of the relations
$q_{n_i}\equiv\pm q_{n_j}$ ($i\neq j$) and $q_{n_i}q_{n_j}\equiv\pm1$
(the orbit collisions), $q_{n_i}^2\equiv-1$ (the mirror condition,
i.e.\ $q_{2n_i}+1\equiv0$) and $q_{n_i}\equiv\pm(2J\pm1)$ (family (b)
of Theorem~\ref{thm:classification}) says that $p$ divides one of the
integers
\[
 a_{n_i}\mp a_{n_j}+2(P_{n_i}\mp P_{n_j})J,\qquad
 a_{n_i+n_j}\mp1+2P_{n_i+n_j}J,\qquad
 a_{2n_i}+1+2P_{2n_i}J,\qquad
 a_{n_i}\mp1+2(P_{n_i}\mp1)J,
\]
each an affine function of $J$ whose coefficient of $J$ is nonzero
($P_n$ is strictly increasing and $P_{n_i}\geq5$).  Finitely many such
functions occur; each is nonzero and of absolute value less than
$p=2J^2-1$ for all large $J$, so $p$ divides none of them.  Family (c)
of Theorem~\ref{thm:classification} is impossible because $p>2J$.
Hence, for all sufficiently large multiples $J$ of $\Lambda$, the $K$
orbits are pairwise distinct, none is mirror, and by
Theorem~\ref{thm:classification} none of the $L(p,q_{n_i})$ is rank-$J$
balanced.
\end{proof}

This deduction from Theorem~\ref{thm:P} and
Theorem~\ref{thm:classification} says that the single
number $\eta_J$ can lose an unbounded amount of information about the
handedness of the geometry: at a suitable rank, as many pairwise
non-homeomorphic chiral lens spaces as desired share the reading
$\eta_J=0$ while having unequal level counts.  It says nothing about the
full level sequences $(D_J(n))_{n\geq0}$, the all-rank sequences
$(\eta_J)_{J\geq1}$ or other invariants, which may well separate them.

\begin{remark}[Rank one against higher rank]
\label{rem:rank-one}
At $J=1$ the opposite holds: $\eta_1(L(p,q))=-4s(q,p)$ vanishes iff
$q^2\equiv-1\pmod p$ \cite[p.~36]{RademacherGrosswald1972}, iff $L(p,q)$
is amphicheiral, and then it is balanced at every level
\cite[Cor.~4.4]{LensChirality2026}; so a zero of $\eta_1$ is exactly a
mirror symmetry.  (As an independent check of the convention,
$\eta_1=0\iff q^2\equiv-1$ was re-verified with the literal Dedekind sum
for every $2\leq p<400$ and every unit $q$: $191$ zeros among $48{,}517$
pairs, no exception; Appendix~\ref{app:verification}.)  At prime moduli
$p\geq5$ the dichotomy survives at every rank with $2J^2\not\equiv1\pmod p$
and $p\nmid J$ (Corollary~\ref{cor:E}(c)); at the blind ranks $p\mid J$
every $q\not\equiv\pm1$ is a zero \cite[Cor.~6.4(d)]{LensChirality2026}.  Corollary~\ref{cor:ambiguity} shows that at
silver ranks it fails without bound.  In the language of chemistry, the
chiral spaces of Corollary~\ref{cor:ambiguity} are cryptochiral for
$\eta_J$: handed, with a vanishing net helicity.
\end{remark}

\section{Open problems}
\label{sec:open}

\begin{enumerate}[label=(\arabic*)]
\item \emph{The period law beyond norm $\pm2$.}  Prove
Conjecture~\ref{conj:PL} for the orbits with $|m|\geq2$, for example the
fourth family $2\delta=(1,2)$ and the orbit $(1,-3)$ of $q\equiv4-4J$,
with effective onsets.
\item \emph{Symmetries.}  The reflection $j\mapsto-j$ of the zero classes
of $\delta^n$ and the mirror relation between the two signs of the
norm-two orbits are observed only.
\item \emph{Zeros outside the $S$-units.}  At $J=20,30,40$ there are about
$6J$ zeros per rank, most of them with an empty box
(Observation~\ref{obs:window}(b)); they are not classified.
\item \emph{Divisors.}  At divisors $p<2J^2-1$ of $2J^2-1$ only the
universal orbits of Observation~\ref{obs:rho-statistics}(a) are proved to
be zeros; Sections~\ref{sec:period-law} and \ref{sec:period-law-proved}
concern the level $p=2J^2-1$ only.
\item \emph{A finite description.}  Whether the full zero set at
$p=2J^2-1$ admits a finite parametrized description is open
(Observation~\ref{obs:window}(c)).
\item \emph{Literature.}  The Dedekind-sum evaluations at Pell-type
arguments used here, and Corollary~\ref{cor:S}, are likely classical;
no literature search has been made for them.
\end{enumerate}

\appendix

\section{Exact computations}
\label{app:verification}

Every statement was accompanied by exact computations in integer or
rational arithmetic, several with predictions fixed before the
corresponding rows were computed (``held-out''; ``exploratory'' rows were
available when the statement was formulated).  They test the identities
and conventions; they do not replace the proofs.  The independent
evaluator named below is the pinned exact implementation of the master
formula used throughout the companion \cite{LensChirality2026}.

\begin{small}
\begin{longtable}{@{}L{0.24\textwidth}L{0.41\textwidth}L{0.2\textwidth}L{0.07\textwidth}@{}}
\toprule
statement & what was checked & rows & failures\\
\midrule
\endhead
Remark~\ref{rem:rank-one} & $s(q,p)=0\iff q^2\equiv-1\pmod p$ by the
literal Dedekind sum, every $2\leq p<400$, every unit $q$
(direct Dedekind sums) & $48{,}517$ ($191$
zeros) & $0$\\
Observation~\ref{obs:rho-statistics}(a), Definition~\ref{def:cayley} &
re-run of the exploration searches, ranges recovered from the retained
scripts and the outcome predicted before computation: the $630$ $S$-units
$s\cdot2^c\sqrt2^{\,a}\delta^b$ ($|c|\leq2$, $|a|\leq3$, $|b|\leq4$) and
$\rho=\pm\delta^k$ ($k\leq14$) at the $57$ silver primes
$17\leq p\leq673$, pinned exact evaluator; universal sets exactly as
predicted; the named orbits lie in $D_p$ for every divisor $p>1$ of
$2J^2-1$, $J\leq100$; counterexamples at $p=49$ &
$57$ primes; $1{,}734$ values & $0$\\
Theorem~\ref{thm:C4} (conjectural stage) & the value
$(J-1)\ind_{J\text{ odd}}$ predicted before computation on every
$J\in[19,60]$ with $p=2J^2-1$ prime; the anchors $p=17,31,97,127$
re-verified; distinctness of the orbit of $3-4J$ from the four orbits of
Lemma~\ref{lem:sunit-four}, family (b), the orbit of
Theorem~\ref{thm:C3} and the mirror orbit on every row; the inverse
orbit $3+4J$ gives the same value on every row & $19$ (held-out) & $0$\\
Conjecture~\ref{conj:C5} & $\eta_J(L(p,3-4J))\neq0$ predicted for every
prime $p\equiv1\pmod8$ with $761\leq p\leq1553$ at its minimal silver
rank & $28$ (held-out) & $1$ ($p=1129$, $J=203$: $\eta=0$)\\
Theorem~\ref{thm:C4} & each of $s$, $R_J$, $\Delta_J$, $L_J$ against its
closed form (direct Dedekind sum, direct residues, enumeration of the
triangle, enumeration of $W\times W$), the assembled value against
\eqref{eq:C4}, and against the independent evaluator, on every $J\leq60$
with $p$ prime ($31$ exploratory rows), every $2\leq J\leq80$ with $p$
composite ($41$) and every $61\leq J\leq100$ with $p$ prime ($14$);
controls: dropping $L_J$ changes every odd-$J$ row; the inverse orbit
gives the same value ($18$ rows); the control premise ``$q=4-4J$ is
nonzero on some even-$J$ row'' failed ($\eta=0$ on all $8$ even rows
checked, Remark~\ref{rem:C4-history}(iii)) & $86$ ($55$ held-out, $41$
composite) & $0$ (one control premise failed)\\
Theorem~\ref{thm:P}, Lemma~\ref{lem:S} & on every row: the unit identity
$q_n(-\nu_nt)\equiv1$; $s(q_n,p)$ by a generic reciprocity algorithm
against \eqref{eq:P-i-dedekind}, \eqref{eq:P-ii-dedekind} (against the
literal Dedekind sum for $J\leq120$; against Corollary~\ref{cor:S} on
every $n=2$ row); Lemma~\ref{lem:S}; $R_J$, $\Delta_J$, $L_J$ from their
definitions against \eqref{eq:P-i-residue}--\eqref{eq:P-ii-lowmode}
(against the literal functions of the previous row for $J\leq100$);
$\eta_J$ against the claim (against the pinned evaluator for $J\leq60$).
Rows: slice (i) $n\leq10$, $k\leq6$ ($59$, exploratory) and $42$ rows never
computed before with $n\leq16$; slice (ii) $n\leq10$, $k\leq8$ ($72$,
exploratory) and $39$ never computed before with $n\leq16$; largest
$J=941{,}664$ (slice (i)) and $941{,}665$ (slice (ii)).  Faulhaber-sum
identities (Remark~\ref{rem:faulhaber-certificate}): four identity sets
by sympy and by Wolfram Engine~14.3.  Fast evaluator against the literal
terms and the pinned evaluator: $106$ rows.  Controls: $R_J-4Js\neq\eta_J$
on all $39$ held-out slice-(ii) rows; $\eta(q^*)=\eta(q)$,
$\eta(-q)=-\eta(q)$ on the $65$ held-out rows with $n\leq12$; the mutated
modulus $a_n$ refuted at $n=3,5,7$ ($9$ fresh rows) & $212$ ($81$
held-out); $8$; $106$; $9$ & $0$\\
Conjecture~\ref{conj:PL} & one polynomial per class modulo $|y|$ fitted
on $J\in[150,400]$ for each of the $64$ orbits of $\mathcal W$ (no fit
failure); $12{,}800$ values at $J\in[1401,1600]$ predicted from the fits
before computation (summary counts retained, not the individual
values); on these ranges the identically vanishing orbits are exactly
$(0,1)$, $(1,-1)$, $(1,0)$, $(1,2)$ & $12{,}800$ (held-out) & $0$\\
Corollary~\ref{cor:U}(a) & one exact nonzero value $\eta_J$ at a rank
$J\leq6$ for each of the $60$ orbits of $\mathcal W$ outside the four
families, each cross-checked against the pinned independent evaluator
(ledger printed in Section~\ref{sec:period-law}) & $60$ & $0$\\
Lemma~\ref{lem:box} & the box sum against the literal $\Delta_J$ and
$L_J$ on random $(p,q,J)$ with $2\leq p\leq150$, $J\leq40$ (draws with
$\gcd(q,p)>1$ skipped); control: the weight $J-|u|-|v|$ fails & $181$
(held-out) & $0$\\
Theorem~\ref{thm:Q} & on every row: $q\cdot2\nu_nt\equiv1$, empty box,
closed form of $s$, $\eta_J$ against the claim; $n\leq8$, $2\leq k\leq6$,
both signs (exploratory), and $n\leq12$ with $k\in\{7,8,9\}$, $13\leq n\leq16$
with $k\in\{1,2\}$, both signs (held-out; largest $J=1{,}331{,}714$);
control: the sign-flipped claim fails on every odd-$n$ held-out row & $70$;
$82$ (held-out) & $0$\\
Lemmas~\ref{lem:N}, \ref{lem:N0} & $J=401,\ldots,410$ and $10{,}000$
(Lemma~\ref{lem:N}); $J=401,\ldots,410$ and $5{,}000$
(Lemma~\ref{lem:N0}) & $22$ (held-out) & $0$\\
Theorem~\ref{thm:PL} (decision procedure) & thresholds and per-class
polynomials for the $24$ families $\delta^n$, $n\leq10$, and
$a_nJ\pm P_n$, $n\leq8$ ($\delta^{10}$ and $n=8$ never computed
before); every class predicted at $k^*_j+3$, $k^*_j+4$ (and $3k^*_j+11$
for $\delta^n$, $n\leq8$, and $a_nJ\pm P_n$, $n\leq6$) before
computation; control: period $P_n+1$ mispredicts $6,13,30,71,170$
classes for $n=3,\ldots,7$; symbolic certificates by sympy and by
Wolfram Engine~14.3 (Remark~\ref{rem:period-certificates}) & $13{,}080$
(held-out); $7$ & $0$\\
Theorem~\ref{thm:PL} (complete zero sets) & for each class of the $24$
families: the class polynomial from three exact values at the
threshold; its integer roots past the threshold (exact integer square
roots, with a pinned self-test on a root of size $10^{18}$), each
re-evaluated exactly; every rank below the threshold evaluated exactly;
the active breakpoints $\leq J+1$ on every class $j<40$, $k\leq5$; box
parity.  Result: one isolated zero ($J=429$), $363$ zeros below the
thresholds, every class polynomial of degree $\leq1$ & $24$ families
& $0$\\
Theorem~\ref{thm:PL} (degree $\leq1$) & the $k^2$ coefficient of each of
the four parts of the bulk (Dedekind, constant, layer cake,
$\alpha$-part) as $\lim(\text{part})/k^2$ from exact expressions, in
$\mathsf S_0$ ($\nu=\pm1$) and $\mathsf S_1$ ($\nu,\sigma=\pm1$), by sympy
and by Wolfram Engine~14.3; corroboration: the
$6{,}025$ retained class triples of the decision-procedure run all have
degree $\leq1$ & $6$ cases; $6{,}025$ & $0$\\
\bottomrule
\end{longtable}
\end{small}

Negative controls (a wrong sign, a dropped term, a wrong modulus, a wrong
period) were evaluated in each block and failed at the stated
counterexamples.  Two nonvanishing predictions were refuted and are
recorded above: Conjecture~\ref{conj:C5} at $p=1129$, and the control
premise of Theorem~\ref{thm:C4} at $q=4-4J$.  The supplementary
integer-root routine uses exact integer arithmetic and includes a
regression test above $10^{18}$; all retained class polynomials here are
linear.

\paragraph{Scope of the computations.}  The accompanying supplement
\cite{LCrecord} supplies exact evaluators and retained calculation records
for the finite checks summarized above.  These checks cover the stated
parameter ranges; the universal statements rest on the written proofs.
The counterexamples to Conjecture~\ref{conj:C5} and to the
$q=4-4J$ control premise remain part of the reported findings.

\section*{Authorship, assistance and rights}
Luqman Mahmood directed this work and is responsible for its claims and final
text. OpenAI Codex and Anthropic Claude assisted with mathematical development,
computation, checking and editing. Internal and automated checks are not
external peer review. Computational tests do not replace proofs of universal
statements, and no physical identification or experiment is implied.
The paper and mathematical data are licensed CC BY 4.0. Software source and
executable mathematical verification scripts are excluded from that grant;
their rights are reserved unless explicitly stated otherwise.

\begin{thebibliography}{99}

\bibitem{MahmoodP1Collection2026}
L. Mahmood, \emph{Eta invariants of transverse-traceless tensor curl on
three-dimensional lens spaces}, companion preprint, 2026,
\texttt{doi:10.5281/zenodo.21433675}.

\bibitem{LensChirality2026}
L. Mahmood, \emph{Spectral chirality of lens spaces: level-wise helicity
balance, eta congruences, and silver zeros for transverse-traceless
tensor curl}, companion preprint, 2026. DOI:
\url{https://doi.org/10.5281/zenodo.23049613}.

\bibitem{BGW2017}
T. Bogart, J. Goodrick and K. Woods, Parametric Presburger arithmetic:
logic, combinatorics, and quasi-polynomial behavior, \emph{Discrete
Analysis} \textbf{2017}:4, 34 pp., \texttt{doi:10.19086/da.1254}.

\bibitem{Boldt2017}
S. Boldt, Properties of the Dirac spectrum on three dimensional lens
spaces, \emph{Osaka J. Math.} \textbf{54} (2017), 747--765; arXiv:1504.03121.

\bibitem{PrzytyckiYasuhara2003}
J. H. Przytycki and A. Yasuhara, Symmetry of links and classification of
lens spaces, \emph{Geom. Dedicata} \textbf{98} (2003), 57--61;
arXiv:math/0011119.

\bibitem{RademacherGrosswald1972}
H. Rademacher and E. Grosswald, \emph{Dedekind sums}, Carus Mathematical
Monographs \textbf{16}, Mathematical Association of America, 1972.

\bibitem{LCrecord}
L. Mahmood, \emph{Silver-rank zeros: computational supplement}, 2026,
distributed with this paper; see \path{supplement/README.md}.
DOI: \url{https://doi.org/10.5281/zenodo.23049615}.

\end{thebibliography}

\end{document}
